\section{Kısmi Türevler, Teğet Düzlemler ve Lineer Yaklaşımlar}

\subsection{Amaç ve sezgi}

Tek değişkenli analizde $y=f(x)$ fonksiyonunun türevi $f'(x)$, $x$ biraz değiştiğinde $y$'nin \emph{ne hızla} değiştiğini
ölçer; geometrik olarak da grafiğe çizilen teğet doğrusunun eğimidir. Mühendislikte ise ilgilendiğimiz büyüklükler çoğu
zaman birden fazla değişkene bağlıdır: bir silindirin hacmi $V=\pi r^2h$ hem yarıçapa hem yüksekliğe, bir metal levhadaki
sıcaklık $T(x,y)$ iki konum koordinatına, ideal bir gazın basıncı $P=nRT/V$ sıcaklığa ve hacme bağlıdır.

Bu durumda ``türev'' sorusu şu hâle gelir: \emph{Diğer değişkenleri sabit tutup yalnızca birini değiştirirsek fonksiyon
ne hızla değişir?} Bu sorunun cevabı \vurg{kısmi türevdir}. Bu konuda sırasıyla şunları öğreneceğiz:
\begin{itemize}
  \item Kısmi türevin limit ile tanımı, hesaplanması ve geometrik anlamı (yüzey üzerindeki izlerin eğimleri).
  \item Yüksek mertebeden ve karma kısmi türevler; karma türevlerin hangi koşulda eşit olduğunu söyleyen
        \textbf{Schwarz (Clairaut) teoremi}.
  \item İki değişkenli bir fonksiyonun \textbf{diferansiyellenebilir} olmasının ne demek olduğu (artım formülü
        $\Delta z=f_x\Delta x+f_y\Delta y+\varepsilon_1\Delta x+\varepsilon_2\Delta y$).
  \item Bir yüzeye bir noktada \textbf{teğet düzlem}, fonksiyonun \textbf{lineerleştirmesi} $L(x,y)$ ve bu yaklaşımın
        \textbf{hata sınırı}.
  \item \textbf{Tam (toplam) diferansiyel} $dz=f_x\,dx+f_y\,dy$, ikinci ve yüksek mertebeden diferansiyeller, ve
        diferansiyelle \textbf{yaklaşık hesap} (mutlak, bağıl ve yüzde değişim).
  \item Bütün bunların üç değişkenli $w=f(x,y,z)$ fonksiyonlarına genellenmesi.
\end{itemize}

\begin{bilgi}[Ana fikir: ``yerel olarak düzlem'']
Tek değişkende, türevlenebilir bir eğriye yeterince yakından bakınca eğri neredeyse bir doğru (teğet doğrusu) gibi
görünür. İki değişkende de aynısı geçerlidir: diferansiyellenebilir bir yüzeye bir noktanın çok yakınından bakınca yüzey
neredeyse bir \emph{düzlem} (teğet düzlem) gibi görünür. Lineerleştirme, tam diferansiyel ve yaklaşık hesap bu tek fikrin
farklı yüzleridir.
\end{bilgi}

\begin{bilgi}[Bu konunun ders kaynakları]
Bu bölüm 22.10 (defter notu), 05.11 ve 08.11 derslerinin tahta notlarına dayanır; Kreyszig (\emph{Advanced Engineering
Mathematics}), Zill \& Wright (\emph{Calculus}, çev. Cangül) ve Balcı (\emph{Genel Matematik II}) kitaplarındaki
anlatımla tamamlanmıştır. Zincir kuralı, kapalı fonksiyonların türevi, yönlü türev ve gradyan \textbf{Konu 6}'da
anlatılmaktadır; burada yalnızca gerektiğinde ``bkz. Konu 6'' diye atıf yapılmıştır.
\end{bilgi}

%-------------------------------------------------------------------------------------------
\subsection{Kısmi türevin tanımı, hesaplanması ve geometrik yorumu}

\subsubsection{Tanım}

\begin{tahta}[\foto{IMG-20261022-WA0054}, \foto{IMG-20261108-WA0090}]
\textbf{Kısmi Türev:} $A\subseteq\mathbb R^2$ ve $f:A\to\mathbb R$, $z=f(x,y)$ ve $(a,b)\in A$ olsun. Eğer
\[ \lim_{h\to0}\frac{f(a+h,b)-f(a,b)}{h} \]
limiti mevcutsa bu limit değerine $f$ fonksiyonunun $(a,b)$ noktasında \textbf{$x$'e göre kısmi türevi} denir;
$f_x(a,b)$ veya $\left.\dfrac{\partial f}{\partial x}\right|_{(a,b)}$ şeklinde gösterilir.
\end{tahta}

Tahtada yalnızca $x$'e göre kısmi türev yazıldı; $y$'ye göre kısmi türev tamamen benzer şekilde tanımlanır:

\begin{tanim}[$x$ ve $y$'ye göre kısmi türevler]
$z=f(x,y)$ fonksiyonunun $(a,b)$ noktasındaki kısmi türevleri (limitler mevcut olmak koşuluyla)
\[
f_x(a,b)=\lim_{h\to0}\frac{f(a+h,b)-f(a,b)}{h},\qquad
f_y(a,b)=\lim_{k\to0}\frac{f(a,b+k)-f(a,b)}{k}.
\]
Bu limitleri her $(x,y)$ noktasında düşünürsek $f_x(x,y)$ ve $f_y(x,y)$ yeni birer iki değişkenli fonksiyon olur.
\end{tanim}

\textbf{Gösterimler.} $z=f(x,y)$ için aşağıdaki yazılışların hepsi aynı anlamdadır:
\[
f_x=\frac{\partial f}{\partial x}=\frac{\partial z}{\partial x}=z_x=D_1f,\qquad
f_y=\frac{\partial f}{\partial y}=\frac{\partial z}{\partial y}=z_y=D_2f .
\]
$\partial$ sembolü (``kısmi'' ya da ``del'' diye okunur), türevin \emph{diğer değişkenler sabitken} alındığını
hatırlatır. Tek değişkenli türevdeki $d$ yerine bilerek farklı bir sembol kullanılır.

\begin{bilgi}[Tanımın anlamı]
$f_x(a,b)$'yi hesaplarken $y$ değişkenini $b$ değerinde \emph{dondururuz}. O zaman $g(x)=f(x,b)$ tek değişkenli bir
fonksiyondur ve $f_x(a,b)=g'(a)$ olur. Benzer şekilde $G(y)=f(a,y)$ ise $f_y(a,b)=G'(b)$. Yani kısmi türev aslında
\emph{tek değişkenli türevdir}; bütün bildiğimiz türev kuralları (toplam, çarpım, bölüm, zincir) aynen geçerlidir.
\end{bilgi}

\subsubsection{Hesaplama kuralı}

\begin{teorem}[Kısmi türev nasıl hesaplanır?]
\begin{itemize}
  \item $f_x$'i bulmak için $y$'yi \textbf{sabit} kabul et ve $f(x,y)$'nin $x$'e göre olağan türevini al.
  \item $f_y$'yi bulmak için $x$'i \textbf{sabit} kabul et ve $f(x,y)$'nin $y$'ye göre olağan türevini al.
\end{itemize}
Sabit kabul edilen değişken, türev alırken tıpkı bir sayı (örneğin $5$) gibi davranır.
\end{teorem}

\begin{soru}[Örnek 1 \foto{IMG-20261022-WA0055}, \foto{IMG-20261108-WA0090}]
$f(x,y)=x^2y^3+3x$ fonksiyonunun birinci mertebeden kısmi türevlerini bulunuz.
\end{soru}
\begin{cevap}
$f_x$: $y$ sabit. $x^2y^3$ teriminde $y^3$ sabit bir katsayıdır; $(x^2)'=2x$ olduğundan türevi $2xy^3$; $3x$'in türevi $3$:
\[ f_x=2xy^3+3 . \]
$f_y$: $x$ sabit. $x^2y^3$ teriminde $x^2$ katsayıdır, $(y^3)'=3y^2$; $3x$ ise $y$'ye göre sabit olduğundan türevi $0$:
\[ f_y=3x^2y^2 . \]
\end{cevap}

\begin{soru}[\ekal Örnek 2 (tanımla hesap)]
$f(x,y)=x^2+3xy$ için $f_x(1,2)$'yi doğrudan limit tanımından hesaplayınız ve kuralla bulunan sonuçla karşılaştırınız.
\end{soru}
\begin{cevap}
$f(1,2)=1+6=7$ ve $f(1+h,2)=(1+h)^2+6(1+h)=7+8h+h^2$. Buna göre
\[
f_x(1,2)=\lim_{h\to0}\frac{f(1+h,2)-f(1,2)}{h}=\lim_{h\to0}\frac{8h+h^2}{h}=\lim_{h\to0}(8+h)=8 .
\]
Kuralla: $f_x=2x+3y$, $f_x(1,2)=2+6=8$. İki yol aynı sonucu verir.
\end{cevap}

\begin{soru}[\ekal Örnek 3 (çeşitli fonksiyonlar)]
Aşağıdaki fonksiyonların birinci mertebeden kısmi türevlerini bulunuz.\\
(a) $f(x,y)=x\,e^{xy}$ \qquad (b) $f(x,y)=\arctan\dfrac{y}{x}$ \qquad (c) $f(x,y)=\sin\!\left(\dfrac{x}{1+y}\right)$
\qquad (d) $f(x,y,z)=x^{yz}$ $(x>0)$
\end{soru}
\begin{cevap}
\textbf{(a)} Çarpım kuralı ($y$ sabit): $f_x=e^{xy}+x\cdot y\,e^{xy}=(1+xy)e^{xy}$. $x$ sabitken $x$ bir katsayıdır:
$f_y=x\cdot x\,e^{xy}=x^2e^{xy}$.

\textbf{(b)} $(\arctan u)'=\dfrac{u'}{1+u^2}$. $y$ sabitken $u=y/x$, $u_x=-y/x^2$:
\[ f_x=\frac{-y/x^2}{1+y^2/x^2}=\frac{-y}{x^2+y^2},\qquad f_y=\frac{1/x}{1+y^2/x^2}=\frac{x}{x^2+y^2}. \]

\textbf{(c)} Zincir kuralı (tek değişkenli): $f_x=\cos\!\left(\dfrac{x}{1+y}\right)\cdot\dfrac{1}{1+y}$,\quad
$f_y=\cos\!\left(\dfrac{x}{1+y}\right)\cdot\dfrac{-x}{(1+y)^2}$.

\textbf{(d)} $x$'e göre türevde üs ($yz$) sabittir: $f_x=yz\,x^{yz-1}$. $y$'ye göre türevde taban sabit, üs değişkendir;
$(a^{u})'=a^{u}\ln a\cdot u'$ kuralıyla $f_y=x^{yz}\ln x\cdot z$, benzer şekilde $f_z=x^{yz}\ln x\cdot y$.
\end{cevap}

\subsubsection{Geometrik yorum: izler ve eğimler}

\begin{bilgi}[Kaynaklardan eklendi]
Kısmi türevin geometrik yorumu derste tahtada ayrıca çizilmedi; aşağıdaki anlatım kaynak kitaplara göre eklenmiştir.
\end{bilgi}

$z=f(x,y)$'nin grafiği uzayda bir $S$ yüzeyidir ve $P(a,b,c)$, $c=f(a,b)$, bu yüzey üzerinde bir noktadır.
\begin{itemize}
  \item $y=b$ düşey düzlemi $S$ yüzeyini bir $C_1$ eğrisi boyunca keser (bu eğriye yüzeyin $y=b$ düzlemindeki
        \textbf{izi} denir). $C_1$'in denklemi $z=f(x,b)$'dir; yani $C_1$, $g(x)=f(x,b)$ fonksiyonunun grafiğidir.
        Dolayısıyla $f_x(a,b)=g'(a)$, \textbf{$C_1$ izinin $P$ noktasındaki teğetinin eğimidir}.
  \item Benzer şekilde $x=a$ düzlemi yüzeyi $C_2$ izi ($z=f(a,y)$) boyunca keser ve $f_y(a,b)$, $C_2$'nin $P$'deki
        teğetinin eğimidir.
\end{itemize}
Fiziksel olarak: $f_x(a,b)$, $y$ sabit tutulup $x$ artırıldığında $z$'nin anlık değişim hızı; $f_y(a,b)$ ise $x$ sabit
tutulup $y$ artırıldığında $z$'nin anlık değişim hızıdır. Yüzeyi bir dağ olarak düşünürsek $f_x$ doğu yönünde, $f_y$
kuzey yönünde yürürken karşılaştığımız yokuşun dikliğidir.

\begin{center}
\begin{tikzpicture}
\begin{axis}[view={115}{28},width=10cm,height=8cm,xlabel=$x$,ylabel=$y$,zlabel=$z$,
  xmin=-1.6,xmax=1.8,ymin=-1.3,ymax=1.5,zmin=-2.5,zmax=4.5,ticks=none]
\addplot3[mesh,colormap={gr}{color(0cm)=(gri!40); color(1cm)=(gri!40)},domain=-1.5:1.5,y domain=-1.2:1.2,samples=17,samples y=13]{4-x^2-2*y^2};
\addplot3[very thick,vurgu,domain=-1.5:1.5,samples=25,samples y=0]({x},{1},{2-x^2});
\addplot3[very thick,yesil,domain=-1.2:1.2,samples=25,samples y=0]({1},{x},{3-2*x^2});
\addplot3[thick,vurgu,dashed,domain=0.3:1.7,samples=2,samples y=0]({x},{1},{1-2*(x-1)});
\addplot3[thick,yesil,dashed,domain=0.65:1.35,samples=2,samples y=0]({1},{x},{1-4*(x-1)});
\addplot3[only marks,mark=*,mark size=2pt,black] coordinates {(1,1,1)};
\node[above left,fill=white,inner sep=1pt] at (axis cs:1,1,1.2) {$P(1,1,1)$};
\node[vurgu,fill=white,inner sep=1pt] at (axis cs:-1.5,1,-2.3) {$C_1$: $y=1$};
\node[yesil,fill=white,inner sep=1pt] at (axis cs:1,-1.2,-0.9) {$C_2$: $x=1$};
\end{axis}
\end{tikzpicture}
\end{center}
\begin{center}
{\small Şekil 1. $z=4-x^2-2y^2$ yüzeyi; $y=1$ düzlemindeki iz $C_1$ (kırmızı) ve $x=1$ düzlemindeki iz $C_2$ (yeşil).
Kesikli doğrular izlerin $P$'deki teğetleridir; eğimleri sırasıyla $f_x(1,1)=-2$ ve $f_y(1,1)=-4$.}
\end{center}

\begin{soru}[\ekal Örnek 4 (izlerin eğimi)]
$f(x,y)=4-x^2-2y^2$ için $f_x(1,1)$ ve $f_y(1,1)$'i bulup geometrik olarak yorumlayınız (Şekil 1).
\end{soru}
\begin{cevap}
$f_x=-2x$, $f_y=-4y$; dolayısıyla $f_x(1,1)=-2$, $f_y(1,1)=-4$, $f(1,1)=1$.
\begin{itemize}
  \item $y=1$ düzlemindeki iz: $z=4-x^2-2=2-x^2$. Bu parabolün $x=1$'deki eğimi $-2x\big|_{x=1}=-2=f_x(1,1)$.
        Teğet doğrusu: $x=1+t,\ y=1,\ z=1-2t$.
  \item $x=1$ düzlemindeki iz: $z=3-2y^2$; $y=1$'deki eğimi $-4y\big|_{y=1}=-4=f_y(1,1)$.
        Teğet doğrusu: $x=1,\ y=1+t,\ z=1-4t$.
\end{itemize}
Her iki eğim negatiftir: $P$ noktasından $x$ ya da $y$ artacak şekilde ilerlersek yüzeyden aşağı ineriz; $y$ yönündeki
iniş daha diktir.
\end{cevap}

\begin{soru}[\ekal Örnek 5 (değişim hızı yorumu: ideal gaz)]
İdeal gaz yasası $PV=nRT$ ($n,R$ sabit). $P$'yi $V$ ve $T$'nin, $V$'yi $T$ ve $P$'nin, $T$'yi $P$ ve $V$'nin fonksiyonu
olarak düşünüp
\[ \frac{\partial P}{\partial V}\cdot\frac{\partial V}{\partial T}\cdot\frac{\partial T}{\partial P}=-1 \]
olduğunu gösteriniz.
\end{soru}
\begin{cevap}
$P=\dfrac{nRT}{V}\Rightarrow \dfrac{\partial P}{\partial V}=-\dfrac{nRT}{V^2}$;\quad
$V=\dfrac{nRT}{P}\Rightarrow \dfrac{\partial V}{\partial T}=\dfrac{nR}{P}$;\quad
$T=\dfrac{PV}{nR}\Rightarrow \dfrac{\partial T}{\partial P}=\dfrac{V}{nR}$. Çarparsak
\[
-\frac{nRT}{V^2}\cdot\frac{nR}{P}\cdot\frac{V}{nR}=-\frac{nRT}{PV}=-\frac{PV}{PV}=-1 .
\]
Bu örnek, $\partial P/\partial V$ gibi ifadelerin \emph{kesir olmadığını} gösterir: ``sadeleştirseydik'' $+1$ bulurduk.
\end{cevap}

%-------------------------------------------------------------------------------------------
\subsection{Yüksek mertebeden ve karma kısmi türevler; Schwarz--Clairaut teoremi}

$f_x$ ve $f_y$ de iki değişkenli fonksiyonlar olduğundan onların da kısmi türevlerini alabiliriz. Böylece dört tane
\textbf{ikinci mertebeden kısmi türev} elde ederiz.

\begin{tahta}[\foto{IMG-20261022-WA0055}, \foto{IMG-20261108-WA0090}]
\textbf{Yüksek Mertebeden Türevler:} $f_x$, $f_y$; $x$ ve $y$'ye göre türevlenebilen fonksiyonlar olsun.
\begin{align*}
f_{xx}&=(f_x)_x=\frac{\partial}{\partial x}\Big(\frac{\partial f}{\partial x}\Big)=\frac{\partial^2f}{\partial x^2}, &
f_{xy}&=(f_x)_y=\frac{\partial}{\partial y}\Big(\frac{\partial f}{\partial x}\Big)=\frac{\partial^2f}{\partial y\,\partial x},\\
f_{yx}&=(f_y)_x=\frac{\partial}{\partial x}\Big(\frac{\partial f}{\partial y}\Big)=\frac{\partial^2f}{\partial x\,\partial y}, &
f_{yy}&=(f_y)_y=\frac{\partial}{\partial y}\Big(\frac{\partial f}{\partial y}\Big)=\frac{\partial^2f}{\partial y^2}.
\end{align*}
\end{tahta}

\begin{dikkat}[Dikkat: türev alma sırası]
Alt indis gösteriminde türev \emph{soldan sağa} alınır: $f_{xy}=(f_x)_y$, önce $x$'e sonra $y$'ye göre.
Kesir gösteriminde ise \emph{sağdan sola} okunur: $\dfrac{\partial^2f}{\partial y\,\partial x}$ önce $x$'e (paydada sağdaki),
sonra $y$'ye göre türev demektir. Yani $f_{xy}=\dfrac{\partial^2f}{\partial y\,\partial x}$.
$f_{xy}$ ve $f_{yx}$'e \textbf{karma (çapraz) kısmi türevler} denir.
\end{dikkat}

Üçüncü ve daha yüksek mertebeden türevler de aynı şekilde tanımlanır; örneğin
$f_{xyy}=\big((f_x)_y\big)_y=\dfrac{\partial^3f}{\partial y^2\,\partial x}$. İki değişkenli bir fonksiyonun $2^n$ tane
$n$-inci mertebe kısmi türevi vardır (sıraya göre); ancak aşağıdaki teorem sayesinde bunların çoğu birbirine eşittir.

\begin{soru}[Örnek 6 \foto{IMG-20261022-WA0055}, \foto{IMG-20261022-WA0059}, \foto{IMG-20261108-WA0091}]
$f(x,y)=\dfrac{xy}{x^2+y^2}$ fonksiyonunun ikinci mertebeden kısmi türevlerini hesaplayınız ($(x,y)\ne(0,0)$).
\end{soru}
\begin{cevap}[Çözüm (adım adım)]
Kısaltma için $r^2=x^2+y^2$ yazalım (yalnızca yazım kolaylığı).

\textbf{1) $f_x$:} Bölüm kuralı, $y$ sabit; payın türevi $y$, paydanın türevi $2x$:
\[
f_x=\frac{y(x^2+y^2)-2x\cdot xy}{(x^2+y^2)^2}=\frac{yx^2+y^3-2x^2y}{(x^2+y^2)^2}=\frac{y^3-x^2y}{(x^2+y^2)^2}.
\]
\textbf{2) $f_{xx}$:} $f_x$'in $x$'e göre türevi. Payın türevi $-2xy$, paydanın türevi $2(x^2+y^2)\cdot2x$:
\[
f_{xx}=\frac{-2xy(x^2+y^2)^2-2(x^2+y^2)\cdot2x\,(y^3-x^2y)}{(x^2+y^2)^4}.
\]
Pay ve paydayı $(x^2+y^2)$ ile sadeleştirirsek pay $-2xy(x^2+y^2)-4x(y^3-x^2y)=-2x^3y-2xy^3-4xy^3+4x^3y$ olur:
\[
f_{xx}=\frac{2x^3y-6xy^3}{(x^2+y^2)^3}.
\]
\textbf{3) $f_{xy}$:} $f_x$'in $y$'ye göre türevi. Payın türevi $3y^2-x^2$, paydanın türevi $2(x^2+y^2)\cdot2y$:
\begin{align*}
f_{xy}&=\frac{(3y^2-x^2)(x^2+y^2)^2-2(x^2+y^2)\cdot2y\,(y^3-x^2y)}{(x^2+y^2)^4}\\
&=\frac{3x^2y^2+3y^4-x^4-x^2y^2-4y^4+4x^2y^2}{(x^2+y^2)^3}
=\boxed{\frac{-x^4-y^4+6x^2y^2}{(x^2+y^2)^3}} .
\end{align*}
\textbf{4) $f_y$:} $x$ sabit; payın türevi $x$, paydanın türevi $2y$:
\[
f_y=\frac{x(x^2+y^2)-2y\cdot xy}{(x^2+y^2)^2}=\frac{x^3-xy^2}{(x^2+y^2)^2}.
\]
\textbf{5) $f_{yx}$:} $f_y$'nin $x$'e göre türevi. Payın türevi $3x^2-y^2$:
\begin{align*}
f_{yx}&=\frac{(3x^2-y^2)(x^2+y^2)^2-2(x^2+y^2)\cdot2x\,(x^3-xy^2)}{(x^2+y^2)^4}\\
&=\frac{3x^4+3x^2y^2-x^2y^2-y^4-4x^4+4x^2y^2}{(x^2+y^2)^3}
=\boxed{\frac{-x^4-y^4+6x^2y^2}{(x^2+y^2)^3}} .
\end{align*}
\textbf{6) $f_{yy}$:} $f_y$'nin $y$'ye göre türevi (payın türevi $-2xy$):
\[
f_{yy}=\frac{-2xy(x^2+y^2)-4y(x^3-xy^2)}{(x^2+y^2)^3}=\frac{2xy^3-6x^3y}{(x^2+y^2)^3}.
\]
\textbf{Sonuç:} $f_{xy}(x,y)=f_{yx}(x,y)$; karma türevler eşittir. (Fonksiyon $x\leftrightarrow y$ simetrik olduğundan
$f_y$ ve $f_{yy}$, $f_x$ ve $f_{xx}$'ten $x$ ile $y$'nin yer değiştirmesiyle de elde edilebilirdi.)
\end{cevap}

\begin{bilgi}[Tahtadaki ara adımlar hakkında not]
\foto{IMG-20261022-WA0055} (defter notu) üzerindeki $f_{xy}$ ara satırında ``$-4y^2$'' yazılmıştı; doğrusu
$-4y\cdot y^3=-4y^4$'tür (\foto{IMG-20261108-WA0091}'de de aynı yazım kayması var). Aynı defter notunda $f_{xx}$ ara
satırındaki ilk terim ``$2x^3y$'' olarak görünüyor; doğrusu $-2x^3y$'dir (\foto{IMG-20261108-WA0091}'de doğru yazılmış).
Her iki durumda da kutulanan nihai sonuçlar doğrudur.
\end{bilgi}

\begin{tahta}[\foto{IMG-20261022-WA0059}, \foto{IMG-20261108-WA0092}]
$f_x$, $f_y$ türevlenebilen fonksiyonlar ve türevleri sürekli ise $f_{xy}(x,y)=f_{yx}(x,y)$:
\textbf{Schwarz ya da Clairaut teoremi}.
\end{tahta}

Teoremin kaynaklardaki kesin ifadesi şöyledir:

\begin{teorem}[Schwarz (Clairaut) teoremi]
$f$, $(a,b)$ noktasını içeren bir açık bölgede tanımlı olsun. Eğer $f_{xy}$ ve $f_{yx}$ karma türevleri bu bölgede
\textbf{sürekli} ise
\[ f_{xy}(a,b)=f_{yx}(a,b). \]
Daha genel olarak: gerekli tüm kısmi türevler sürekli olduğunda türev alma \emph{sırası önemli değildir}; örneğin
$f_{xyy}=f_{yxy}=f_{yyx}$.
\end{teorem}

\begin{dikkat}[Tahtada düzeltilen ifade]
\foto{IMG-20261022-WA0059} (defter notu) üzerinde ``$f_{xy}=f_{yx}$, eğer $f(x,y)=f(y,x)$ ise geçerlidir'' yazılmıştı.
\textbf{Bu ifade yanlıştır.} Karma türevlerin eşitliği fonksiyonun $x$ ve $y$'ye göre simetrik olmasına bağlı
\emph{değildir}; belirleyici olan karma türevlerin \emph{sürekliliğidir} (yukarıdaki teorem). Örneğin simetrik olmayan
$f=x^2y^3+3x$ için $f_{xy}=(2xy^3+3)_y=6xy^2$ ve $f_{yx}=(3x^2y^2)_x=6xy^2$, yani yine eşittir. Hemen altında yazılan
Schwarz/Clairaut ifadesi doğru koşulu vermektedir.
\end{dikkat}

\begin{soru}[\ekal Örnek 7 (Clairaut doğrulaması)]
$f(x,y)=x^3y^4+\sin(xy^2)$ için $f_{xy}=f_{yx}$ olduğunu doğrulayınız.
\end{soru}
\begin{cevap}
$f_x=3x^2y^4+y^2\cos(xy^2)$,\quad $f_y=4x^3y^3+2xy\cos(xy^2)$.
\begin{align*}
f_{xy}&=12x^2y^3+2y\cos(xy^2)+y^2\cdot\big(-\sin(xy^2)\big)\cdot2xy=12x^2y^3+2y\cos(xy^2)-2xy^3\sin(xy^2),\\
f_{yx}&=12x^2y^3+2y\cos(xy^2)+2xy\cdot\big(-\sin(xy^2)\big)\cdot y^2=12x^2y^3+2y\cos(xy^2)-2xy^3\sin(xy^2).
\end{align*}
İkisi aynıdır; bu beklenen bir sonuçtur çünkü tüm kısmi türevler her yerde süreklidir.
\end{cevap}

\begin{soru}[\ekal Örnek 8 (sürekli olmayan karma türevler: $f_{xy}\ne f_{yx}$)]
\[
f(x,y)=\begin{cases}\dfrac{xy(x^2-y^2)}{x^2+y^2}, & (x,y)\ne(0,0)\\[4pt] 0, & (x,y)=(0,0)\end{cases}
\]
için $f_{xy}(0,0)=-1$, $f_{yx}(0,0)=1$ olduğunu gösteriniz.
\end{soru}
\begin{cevap}
Önce eksenler üzerinde birinci türevleri tanımdan bulalım. $y\ne0$ sabitken
\[
f_x(0,y)=\lim_{h\to0}\frac{f(h,y)-f(0,y)}{h}=\lim_{h\to0}\frac{y(h^2-y^2)}{h^2+y^2}=\frac{-y^3}{y^2}=-y ,
\]
ve $f_x(0,0)=\lim_{h\to0}\frac{0-0}{h}=0$; yani her $y$ için $f_x(0,y)=-y$. Benzer şekilde $x\ne0$ için
\[
f_y(x,0)=\lim_{k\to0}\frac{f(x,k)-f(x,0)}{k}=\lim_{k\to0}\frac{x(x^2-k^2)}{x^2+k^2}=x ,
\]
ve $f_y(0,0)=0$; yani $f_y(x,0)=x$. Şimdi
\[
f_{xy}(0,0)=\frac{d}{dy}\big(f_x(0,y)\big)\Big|_{y=0}=\frac{d}{dy}(-y)=-1,\qquad
f_{yx}(0,0)=\frac{d}{dx}\big(f_y(x,0)\big)\Big|_{x=0}=\frac{d}{dx}(x)=1 .
\]
Karma türevler eşit değildir. Teoremle çelişki yoktur: bu fonksiyonun karma türevleri $(0,0)$'da sürekli değildir.
(İlginç olarak bu fonksiyon da \emph{simetrik değil}, antisimetriktir: $f(y,x)=-f(x,y)$.)
\end{cevap}

\begin{soru}[\ekal Örnek 9 (kısmi diferansiyel denklemler; Kreyszig tarzı)]
(a) $u=\ln(x^2+y^2)$ ve $u=e^{-x}\cos y$ fonksiyonlarının Laplace denklemi $u_{xx}+u_{yy}=0$'ı sağladığını,\\
(b) $u=\sin(x-ct)$ fonksiyonunun dalga denklemi $u_{tt}=c^2u_{xx}$'i,\\
(c) $u=e^{-kt}\sin x$ fonksiyonunun ısı denklemi $u_t=k\,u_{xx}$'i sağladığını gösteriniz.
\end{soru}
\begin{cevap}
\textbf{(a)} $u=\ln(x^2+y^2)$: $u_x=\dfrac{2x}{x^2+y^2}$,
$u_{xx}=\dfrac{2(x^2+y^2)-2x\cdot2x}{(x^2+y^2)^2}=\dfrac{2y^2-2x^2}{(x^2+y^2)^2}$; simetriden
$u_{yy}=\dfrac{2x^2-2y^2}{(x^2+y^2)^2}$. Toplam $0$. \quad
$u=e^{-x}\cos y$: $u_{xx}=e^{-x}\cos y$, $u_{yy}=-e^{-x}\cos y$; toplam $0$.

\textbf{(b)} $u_t=-c\cos(x-ct)$, $u_{tt}=-c^2\sin(x-ct)$; $u_x=\cos(x-ct)$, $u_{xx}=-\sin(x-ct)$. Dolayısıyla
$u_{tt}=c^2u_{xx}$.

\textbf{(c)} $u_t=-k e^{-kt}\sin x$, $u_{xx}=-e^{-kt}\sin x$; dolayısıyla $u_t=k\,u_{xx}$.
\end{cevap}

%-------------------------------------------------------------------------------------------
\subsection{Diferansiyellenebilirlik ve artım}

\subsubsection{Tek değişkende hatırlatma}

\begin{tahta}[\foto{IMG-20261022-WA0059}, \foto{IMG-20261108-WA0092}]
\textbf{Türevlenebilme -- Tek değişkende:}
\[
\frac{\Delta y}{\Delta x}=\frac{f(x+\Delta x)-f(x)}{\Delta x}=f'(x)+\varepsilon
\ \Rightarrow\ \Delta y=f'(x)\Delta x+\varepsilon\Delta x ;
\]
$\varepsilon$: $\Delta x\to0$ iken sıfıra giden bir büyüklük; $f'(x)\Delta x$: $f$'nin diferansiyeli.
\end{tahta}

Burada önemli olan şudur: $y$'deki gerçek değişim $\Delta y$, $\Delta x$'e göre \emph{lineer} olan $f'(x)\Delta x$ terimi
ile, $\Delta x$'ten \emph{çok daha hızlı} sıfıra giden $\varepsilon\Delta x$ teriminin toplamıdır. Tek değişkende
``türevi var'' ile ``bu biçimde yazılabilir'' aynı şeydir.

\subsubsection{İki değişkende artım ve diferansiyellenebilirlik}

\begin{tahta}[\foto{IMG-20261022-WA0057}, \foto{IMG-20261108-WA0092}]
$f(x,y)$'nin $f_x$ ve $f_y$ kısmi türevleri $(x_0,y_0)$ noktasını içeren bir $B$ bölgesinde tanımlı ve $(x_0,y_0)$
noktasında sürekli olsun. $(x,y)$ noktası $(x_0,y_0)$'dan $(x_0+\Delta x,y_0+\Delta y)$ noktasına kadar değiştiğinde
$z=f(x,y)$ fonksiyonunda meydana gelen değişim
\[ \Delta z=f(x_0+\Delta x,y_0+\Delta y)-f(x_0,y_0) \]
\[ \Rightarrow\ \Delta z=f_x(x_0,y_0)\Delta x+f_y(x_0,y_0)\Delta y+\varepsilon_1\Delta x+\varepsilon_2\Delta y \]
$\varepsilon_1$ ve $\varepsilon_2$: $\Delta x\to0$, $\Delta y\to0$ iken sıfıra giden büyüklüklerdir.

\textbf{Tanım:} $f_x(x_0,y_0)$, $f_y(x_0,y_0)$ kısmi türevleri mevcut ve $\Delta z$ yukarıdaki biçimde yazılabiliyorsa
$f(x,y)$ fonksiyonu \textbf{türevlenebilirdir (diferansiyellenebilirdir)} denir.

\textbf{Sonuç:} $f$, $(x_0,y_0)$ noktasında diferansiyellenebilir ise $(x_0,y_0)$ noktasında süreklidir.
\end{tahta}

Tahtadaki bu not iki farklı şeyi birlikte içerir; kaynaklara göre bunları ayırmak gerekir:

\begin{tanim}[Diferansiyellenebilirlik]
$f_x(x_0,y_0)$ ve $f_y(x_0,y_0)$ mevcut olsun. Eğer $f$'nin artımı
\[
\Delta z=f_x(x_0,y_0)\,\Delta x+f_y(x_0,y_0)\,\Delta y+\varepsilon_1\Delta x+\varepsilon_2\Delta y
\]
biçiminde yazılabiliyor ve $(\Delta x,\Delta y)\to(0,0)$ iken $\varepsilon_1\to0$, $\varepsilon_2\to0$ oluyorsa $f$,
$(x_0,y_0)$ noktasında \textbf{diferansiyellenebilirdir}. Bir bölgenin her noktasında diferansiyellenebilen fonksiyona
o bölgede diferansiyellenebilir denir.
\end{tanim}

\begin{teorem}[Artım teoremi (diferansiyellenebilirlik için yeter koşul)]
$f_x$ ve $f_y$, $(x_0,y_0)$'ı içeren bir açık bölgede tanımlı ve $(x_0,y_0)$'da \textbf{sürekli} ise $f$,
$(x_0,y_0)$'da diferansiyellenebilirdir.
\end{teorem}

\textbf{İspat fikri.} Artımı iki adımda yazalım: önce yalnız $x$'i, sonra yalnız $y$'yi değiştirelim:
\[
\Delta z=\big[f(x_0+\Delta x,y_0+\Delta y)-f(x_0,y_0+\Delta y)\big]+\big[f(x_0,y_0+\Delta y)-f(x_0,y_0)\big].
\]
Her köşeli parantez tek değişkenli bir farktır; ortalama değer teoreminden $x_0$ ile $x_0+\Delta x$ arasında bir $c$ ve
$y_0$ ile $y_0+\Delta y$ arasında bir $d$ için
$\Delta z=f_x(c,\,y_0+\Delta y)\,\Delta x+f_y(x_0,d)\,\Delta y$ olur.
$\varepsilon_1=f_x(c,y_0+\Delta y)-f_x(x_0,y_0)$ ve $\varepsilon_2=f_y(x_0,d)-f_y(x_0,y_0)$ dersek, kısmi türevlerin
sürekliliğinden $\Delta x,\Delta y\to0$ iken $\varepsilon_1,\varepsilon_2\to0$ olur. \hfill$\square$

\textbf{Sonucun gerekçesi.} Diferansiyellenebilirse $\Delta z=(f_x+\varepsilon_1)\Delta x+(f_y+\varepsilon_2)\Delta y$
ifadesi $\Delta x,\Delta y\to0$ iken $0$'a gider; yani $f(x_0+\Delta x,y_0+\Delta y)\to f(x_0,y_0)$, bu da süreklilik
demektir.

\begin{bilgi}[Eşdeğer tanım]
Bazı kitaplarda (ör. Kreyszig) diferansiyellenebilirlik şöyle verilir:
\[
\lim_{(\Delta x,\Delta y)\to(0,0)}
\frac{\Delta z-f_x(x_0,y_0)\Delta x-f_y(x_0,y_0)\Delta y}{\sqrt{\Delta x^2+\Delta y^2}}=0 .
\]
Yani ``gerçek değişim ile lineer kısım arasındaki fark, adım uzunluğundan çok daha hızlı küçülür''. Bu, tahtadaki
$\varepsilon_1,\varepsilon_2$ tanımıyla eşdeğerdir.
\end{bilgi}

Mantıksal ilişkiler aşağıdaki şemada özetlenmiştir. Tek değişkenden farklı olarak, iki değişkende \emph{kısmi türevlerin
var olması, diferansiyellenebilirliği (hatta sürekliliği bile) garanti etmez.}

\begin{center}
\begin{tikzpicture}[kut/.style={draw=ana,thick,rounded corners,fill=acik,align=center,font=\small,
  minimum height=1cm,text width=3.4cm}]
\node[kut] (c1) at (0,0) {$f_x$, $f_y$ var ve\\ sürekli};
\node[kut] (d) at (4.6,0) {$f$ diferansiyel-\\lenebilir};
\node[kut] (s) at (9.2,0.9) {$f$ sürekli};
\node[kut] (k) at (9.2,-0.9) {$f_x$, $f_y$ var};
\draw[-{Stealth[length=3mm]},very thick,yesil] (c1) -- (d);
\draw[-{Stealth[length=3mm]},very thick,yesil] (d.east) -- (s.west);
\draw[-{Stealth[length=3mm]},very thick,yesil] (d.east) -- (k.west);
\draw[-{Stealth[length=3mm]},thick,vurgu,dashed] (k.south) to[bend left=25] node[below,font=\small]{$\times$ genelde doğru değil} (d.south);
\end{tikzpicture}
\end{center}
\begin{center}
{\small Şekil 2. Süreklilik, kısmi türevlerin varlığı ve diferansiyellenebilirlik arasındaki ilişkiler.}
\end{center}

\begin{soru}[\ekal Örnek 10 ($\varepsilon_1$, $\varepsilon_2$'yi açıkça bulma)]
(a) $f(x,y)=x^2+y^2$ ve (b) $f(x,y)=xy$ fonksiyonlarının her noktada diferansiyellenebilir olduğunu, $\varepsilon_1$
ve $\varepsilon_2$'yi açıkça bularak gösteriniz.
\end{soru}
\begin{cevap}
\textbf{(a)} $\Delta z=(x+\Delta x)^2+(y+\Delta y)^2-x^2-y^2=2x\Delta x+2y\Delta y+\Delta x\cdot\Delta x+\Delta y\cdot\Delta y$.
$f_x=2x$, $f_y=2y$ olduğundan $\varepsilon_1=\Delta x$, $\varepsilon_2=\Delta y$ seçilir; ikisi de
$(\Delta x,\Delta y)\to(0,0)$ iken $0$'a gider.

\textbf{(b)} $\Delta z=(x+\Delta x)(y+\Delta y)-xy=y\Delta x+x\Delta y+\Delta x\,\Delta y$. $f_x=y$, $f_y=x$;
$\varepsilon_1=\Delta y$, $\varepsilon_2=0$ alınabilir (ya da $\varepsilon_1=0$, $\varepsilon_2=\Delta x$:
$\varepsilon$'lar tek türlü belirli değildir).
\end{cevap}

\begin{soru}[\ekal Örnek 11 (kısmi türevler var, ama fonksiyon diferansiyellenemez)]
\[
f(x,y)=\begin{cases}\dfrac{xy}{x^2+y^2}, & (x,y)\ne(0,0)\\[4pt] 0, & (x,y)=(0,0)\end{cases}
\]
fonksiyonunun (Örnek 6'daki fonksiyonun orijinde $0$ ile tanımlanmış hâli) $(0,0)$'da kısmi türevlerinin var olduğunu
fakat $f$'nin orada diferansiyellenemez olduğunu gösteriniz.
\end{soru}
\begin{cevap}
Tanımdan: $f(h,0)=0$ olduğundan $f_x(0,0)=\lim_{h\to0}\frac{0-0}{h}=0$; benzer şekilde $f_y(0,0)=0$. Kısmi türevler vardır.

Fakat $y=mx$ doğrusu boyunca $f(x,mx)=\dfrac{mx^2}{x^2+m^2x^2}=\dfrac{m}{1+m^2}$, yani limit $m$'ye bağlıdır
($m=0$ için $0$, $m=1$ için $1/2$). $(0,0)$'da limit yoktur, dolayısıyla $f$ orada sürekli değildir. Diferansiyellenebilir
fonksiyon sürekli olmak zorunda olduğundan $f$, $(0,0)$'da \textbf{diferansiyellenemez}.
\end{cevap}

\begin{soru}[\ekal Örnek 12 (sürekli, kısmi türevleri var, yine de diferansiyellenemez)]
$f(x,y)=\sqrt[3]{x^3+y^3}$ fonksiyonu için $f_x(0,0)$, $f_y(0,0)$'ı bulunuz ve $f$'nin $(0,0)$'da diferansiyellenemez
olduğunu gösteriniz.
\end{soru}
\begin{cevap}
$f_x(0,0)=\lim_{h\to0}\dfrac{\sqrt[3]{h^3}-0}{h}=1$, benzer şekilde $f_y(0,0)=1$. $f$ her yerde süreklidir.
Diferansiyellenebilir olsaydı $\Delta z\approx f_x\Delta x+f_y\Delta y=\Delta x+\Delta y$ olmalıydı ve fark,
adım uzunluğuna bölündüğünde $0$'a gitmeliydi. $\Delta x=\Delta y=t$ alalım: $\Delta z=\sqrt[3]{2t^3}=\sqrt[3]2\,t$,
lineer kısım $2t$. Fark $(\sqrt[3]2-2)t$, adım uzunluğu $\sqrt2|t|$; oranın mutlak değeri
$\dfrac{2-\sqrt[3]2}{\sqrt2}\approx0{,}52\ne0$. Dolayısıyla $f$, $(0,0)$'da diferansiyellenemez. (Teoremle çelişki yok:
$f_x=x^2(x^3+y^3)^{-2/3}$ orijinde sürekli değildir.)
\end{cevap}

%-------------------------------------------------------------------------------------------
\subsection{Teğet düzlem}

\begin{bilgi}[Derste tahtada görülmedi, kaynaklardan eklendi]
Elimizdeki tahta fotoğraflarında teğet düzlem konusu işlenmemiştir (\foto{IMG-20261108-WA0093}'ün sağ kenarında yalnızca
kesik bir ``T\ldots'' başlığı görünüyor). Bu alt başlık müfredatta yer aldığı için Zill \& Wright, Balcı ve Kreyszig'e
göre eksiksiz olarak eklenmiştir. Lineerleştirme (tahtada işlendi) ile teğet düzlem aynı formüldür; aşağıda bu bağlantı
açıkça gösterilmiştir.
\end{bilgi}

\subsubsection{Teğet düzlemin elde edilişi}

$z=f(x,y)$ yüzeyi üzerinde $P(x_0,y_0,z_0)$, $z_0=f(x_0,y_0)$ noktasını alalım ve $f_x$, $f_y$ sürekli olsun.
Bir önceki bölümde gördüğümüz iki izin $P$'deki teğet doğrularını düşünelim:
\begin{itemize}
  \item $y=y_0$ düzlemindeki $C_1$ izinin teğeti: $x$ bir birim artarken $z$, $f_x(x_0,y_0)$ kadar artar. Doğrultu
        vektörü $\mathbf u=\langle1,\,0,\,f_x(x_0,y_0)\rangle$.
  \item $x=x_0$ düzlemindeki $C_2$ izinin teğeti: doğrultu vektörü $\mathbf v=\langle0,\,1,\,f_y(x_0,y_0)\rangle$.
\end{itemize}
\textbf{Teğet düzlem}, bu iki teğet doğruyu içeren düzlemdir. Normal vektörü (Konu 1'deki vektörel çarpımla):
\[
\mathbf n=\mathbf u\times\mathbf v=
\begin{vmatrix}\mathbf i&\mathbf j&\mathbf k\\ 1&0&f_x\\ 0&1&f_y\end{vmatrix}
=\langle -f_x,\,-f_y,\,1\rangle
\quad\text{ya da (}-1\text{ ile çarparak)}\quad \mathbf n=\langle f_x,\,f_y,\,-1\rangle .
\]
Noktası ve normali bilinen düzlemin denklemi (Konu 2) $f_x(x-x_0)+f_y(y-y_0)-(z-z_0)=0$ olur.

\begin{teorem}[Teğet düzlem ve normal doğru]
$f_x$ ve $f_y$, $(x_0,y_0)$'da sürekli olsun. $z=f(x,y)$ yüzeyinin $P(x_0,y_0,z_0)$ noktasındaki \textbf{teğet düzlemi}
\[
\boxed{\,z-z_0=f_x(x_0,y_0)\,(x-x_0)+f_y(x_0,y_0)\,(y-y_0)\,}
\]
ve bu düzleme $P$'de dik olan \textbf{normal doğru}
\[
\frac{x-x_0}{f_x(x_0,y_0)}=\frac{y-y_0}{f_y(x_0,y_0)}=\frac{z-z_0}{-1}
\qquad\big(\text{parametrik: }(x,y,z)=(x_0,y_0,z_0)+t\,\langle f_x,f_y,-1\rangle\big).
\]
\end{teorem}

\begin{center}
\begin{tikzpicture}
\begin{axis}[view={-40}{22},width=10cm,height=8cm,xlabel=$x$,ylabel=$y$,zlabel=$z$,
  xmin=-0.6,xmax=2,ymin=-0.6,ymax=2,zmin=-2,zmax=10,ticks=none]
\addplot3[surf,colormap/cool,opacity=0.75,faceted color=gri!60,domain=-0.5:1.9,y domain=-0.5:1.9,samples=15]{2*x^2+y^2};
\addplot3[surf,shader=flat,colormap={duz}{color(0cm)=(vurgu!35); color(1cm)=(vurgu!35)},faceted color=vurgu,opacity=0.8,domain=0.25:1.75,y domain=0.25:1.75,samples=2]{4*x+2*y-3};
\addplot3[only marks,mark=*,mark size=2pt,black] coordinates {(1,1,3)};
\addplot3[thick,ana,-{Stealth[length=2.5mm]},domain=0:0.55,samples=2,samples y=0]({1+4*x},{1+2*x},{3-x});
\node[above left] at (axis cs:1,1,3.2) {$P(1,1,3)$};
\node[ana] at (axis cs:3.2,2.1,2.4) {$\mathbf n$};
\end{axis}
\end{tikzpicture}
\end{center}
\begin{center}
{\small Şekil 3. $z=2x^2+y^2$ paraboloidi ve $P(1,1,3)$ noktasındaki teğet düzlemi $z=4x+2y-3$ (kırmızı yama);
$\mathbf n=\langle4,2,-1\rangle$ normal vektörü.}
\end{center}

\begin{bilgi}[Teğet düzlem ne zaman vardır?]
Teğet düzlemin anlamlı olması için yalnızca $f_x$ ve $f_y$'nin var olması yetmez; $f$'nin $(x_0,y_0)$'da
\emph{diferansiyellenebilir} olması gerekir (Örnek 11'deki fonksiyonun orijinde kısmi türevleri vardır ama yüzey orada
yırtıktır, teğet düzlemi yoktur). Kısmi türevlerin sürekli olması bunu garanti eder.
Kapalı biçimde $F(x,y,z)=0$ olarak verilen yüzeylerin teğet düzlemi gradyan vektörü $\nabla F$ ile bulunur; bkz. Konu 6.
\end{bilgi}

\begin{teorem}[Adım adım yöntem]
\begin{enumerate}
  \item Noktanın $z$ koordinatını bul: $z_0=f(x_0,y_0)$ (verilmişse yüzey üzerinde olduğunu kontrol et).
  \item $f_x$ ve $f_y$'yi genel olarak hesapla, sonra $(x_0,y_0)$'da \textbf{sayısal} değerlerini bul.
  \item $z-z_0=f_x(x_0,y_0)(x-x_0)+f_y(x_0,y_0)(y-y_0)$ formülüne yerleştir ve sadeleştir.
  \item İstenirse normal doğru: doğrultu vektörü $\langle f_x(x_0,y_0),\,f_y(x_0,y_0),\,-1\rangle$.
\end{enumerate}
\end{teorem}

\begin{soru}[\ekal Örnek 13]
$z=2x^2+y^2$ eliptik paraboloidinin $(1,1,3)$ noktasındaki teğet düzleminin ve normal doğrusunun denklemlerini bulunuz.
\end{soru}
\begin{cevap}
$f(1,1)=2+1=3$ (nokta yüzey üzerinde). $f_x=4x\Rightarrow f_x(1,1)=4$; $f_y=2y\Rightarrow f_y(1,1)=2$.
\[
z-3=4(x-1)+2(y-1)\quad\Longrightarrow\quad z=4x+2y-3 .
\]
Normal doğru: $\mathbf n=\langle4,2,-1\rangle$;\quad $x=1+4t,\ y=1+2t,\ z=3-t$, ya da
$\dfrac{x-1}{4}=\dfrac{y-1}{2}=\dfrac{z-3}{-1}$ (Şekil 3).
\end{cevap}

\begin{soru}[\ekal Örnek 14]
Aşağıdaki yüzeylerin verilen noktadaki teğet düzlemlerini bulunuz.\quad
(a) $z=x\,e^{xy}$, $(1,0,1)$\qquad (b) $z=\ln(x-2y)$, $(3,1,0)$
\end{soru}
\begin{cevap}
\textbf{(a)} $f_x=(1+xy)e^{xy}\Rightarrow f_x(1,0)=1$;\ $f_y=x^2e^{xy}\Rightarrow f_y(1,0)=1$.
$z-1=1\cdot(x-1)+1\cdot(y-0)\Rightarrow z=x+y$.

\textbf{(b)} $f(3,1)=\ln1=0$ $\checkmark$. $f_x=\dfrac{1}{x-2y}\Rightarrow f_x(3,1)=1$;\ $f_y=\dfrac{-2}{x-2y}\Rightarrow f_y(3,1)=-2$.
$z-0=(x-3)-2(y-1)\Rightarrow z=x-2y-1$.
\end{cevap}

\begin{soru}[\ekal Örnek 15 (yatay teğet düzlem, verilen düzleme paralel teğet düzlem)]
(a) $z=x^2+y^2-4x+2y$ yüzeyinin hangi noktasında teğet düzlem yataydır?\\
(b) $z=x^2+2y^2$ yüzeyinin hangi noktasındaki teğet düzlem $2x+8y-z=5$ düzlemine paraleldir? O teğet düzlemi yazınız.
\end{soru}
\begin{cevap}
\textbf{(a)} Teğet düzlem yatay $\iff$ normal $\langle f_x,f_y,-1\rangle$, $\mathbf k$'ya paralel $\iff f_x=f_y=0$.
$f_x=2x-4=0\Rightarrow x=2$; $f_y=2y+2=0\Rightarrow y=-1$; $z=4+1-8-2=-5$. Nokta $(2,-1,-5)$, teğet düzlem $z=-5$.

\textbf{(b)} Paralel düzlemlerin normalleri paraleldir. Verilen düzlemin normali $\langle2,8,-1\rangle$; teğet düzleminki
$\langle2x,4y,-1\rangle$. Üçüncü bileşenler eşit olduğundan $2x=2$, $4y=8$: $x=1$, $y=2$, $z=1+8=9$.
Teğet düzlem: $z-9=2(x-1)+8(y-2)\Rightarrow z=2x+8y-9$ (yani $2x+8y-z=9$; verilen düzleme gerçekten paralel).
\end{cevap}

%-------------------------------------------------------------------------------------------
\subsection{Lineerleştirme ve lineer yaklaşımda hata}

\subsubsection{Lineerleştirme}

\begin{tahta}[\foto{IMG-20261022-WA0057}, \foto{IMG-20261108-WA0093}]
\textbf{Lineerleştirme ve hata:} $f$, $(x_0,y_0)$ noktasında türevlenebilir olsun ve $\Delta x=x-x_0$, $\Delta y=y-y_0$:
\[
f(x,y)=f(x_0,y_0)+f_x(x_0,y_0)(x-x_0)+f_y(x_0,y_0)(y-y_0)+\varepsilon_1\Delta x+\varepsilon_2\Delta y
\]
$\Delta x\to0$ ve $\Delta y\to0$ ise $\varepsilon_1$ ve $\varepsilon_2$ sıfıra giden büyüklükler olduğundan:
\[ f(x,y)\cong f(x_0,y_0)+f_x(x_0,y_0)(x-x_0)+f_y(x_0,y_0)(y-y_0) \]
\[ L(x,y)=f(x_0,y_0)+f_x(x_0,y_0)(x-x_0)+f_y(x_0,y_0)(y-y_0) \]
ifadesine $f$ fonksiyonunun \textbf{lineerleştirmesi} denir.
\end{tahta}

\begin{tanim}[Lineerleştirme (lineer yaklaşım)]
$f$, $(x_0,y_0)$'da diferansiyellenebilir olsun.
\[
L(x,y)=f(x_0,y_0)+f_x(x_0,y_0)(x-x_0)+f_y(x_0,y_0)(y-y_0)
\]
fonksiyonuna $f$'nin $(x_0,y_0)$'daki \textbf{lineerleştirmesi}, $(x_0,y_0)$ civarındaki
\[ f(x,y)\approx L(x,y) \]
yaklaşımına da \textbf{lineer yaklaşım} (teğet düzlem yaklaşımı) denir.
\end{tanim}

\begin{bilgi}[Lineerleştirme = teğet düzlem]
$z=L(x,y)$ denklemi, bir önceki bölümdeki teğet düzlem denkleminin $z$'ye göre çözülmüş hâlidir. Yani lineer yaklaşım,
yüzeyi $P$'nin yakınında teğet düzlemle değiştirmek demektir. Tek değişkendeki $f(x)\approx f(x_0)+f'(x_0)(x-x_0)$
(teğet doğru yaklaşımı) formülünün iki boyutlu karşılığıdır.
\end{bilgi}

\begin{soru}[Örnek 16 \foto{IMG-20261105-WA0037}, \foto{IMG-20261108-WA0085}]
$f(x,y)=e^{2y-x}$ fonksiyonunun $P(0,0)$ ve $Q(1,2)$ noktalarındaki lineerleştirmelerini bulunuz.
\end{soru}
\begin{cevap}
Kısmi türevler (zincir kuralı, iç fonksiyon $2y-x$): $f_x=-e^{2y-x}$, $f_y=2e^{2y-x}$.

\textbf{$P(0,0)$:} $f(0,0)=e^0=1$, $f_x(0,0)=-1$, $f_y(0,0)=2$:
\[ L(x,y)=1-1\cdot(x-0)+2(y-0)=1-x+2y . \]
\textbf{$Q(1,2)$:} $2y-x=3$; $f(1,2)=e^3$, $f_x(1,2)=-e^3$, $f_y(1,2)=2e^3$:
\[ L(x,y)=e^3-e^3(x-1)+2e^3(y-2)=e^3\,(1-x+1+2y-4)=e^3(2y-x-2). \]
Not: Lineerleştirme noktaya bağlıdır; aynı fonksiyonun farklı noktalardaki lineerleştirmeleri farklı düzlemlerdir.
\end{cevap}

\begin{soru}[\ekal Örnek 17 (lineer yaklaşımla sayısal hesap)]
$f(x,y)=\sqrt{x^2+y^2}$'nin $(3,4)$'teki lineerleştirmesini bulunuz ve $\sqrt{(3{,}02)^2+(3{,}97)^2}$'yi yaklaşık
hesaplayınız.
\end{soru}
\begin{cevap}
$f(3,4)=5$; $f_x=\dfrac{x}{\sqrt{x^2+y^2}}\Rightarrow f_x(3,4)=\dfrac35$; $f_y=\dfrac{y}{\sqrt{x^2+y^2}}\Rightarrow
f_y(3,4)=\dfrac45$.
\[ L(x,y)=5+\tfrac35(x-3)+\tfrac45(y-4). \]
$x-3=0{,}02$, $y-4=-0{,}03$:\quad $L(3{,}02;\,3{,}97)=5+0{,}012-0{,}024=4{,}988$.\\
Gerçek değer $\sqrt{24{,}8813}\approx4{,}98812$; yaklaşımın hatası yaklaşık $0{,}0001$.
\end{cevap}

\subsubsection{Lineer yaklaşımda hata}

\begin{tahta}[\foto{IMG-20261022-WA0060}, \foto{IMG-20261108-WA0093}]
$E(x,y)=f(x,y)-L(x,y)$ türevli fonksiyonu: \textbf{hata fonksiyonu}.

$f$, $B=\{(x,y):|x-x_0|<h,\ |y-y_0|<h\}$ dikdörtgensel bölgesinde ikinci türevleri sürekli bir fonksiyon olsun.
$|f_{xx}|,|f_{yy}|,|f_{xy}|$'nin en büyük değeri $M$ ise
\[ |E(x,y)|\le\frac12M\big(|x-x_0|+|y-y_0|\big)^2 \]
ile hesaplanır.
\end{tahta}

\begin{teorem}[Lineer yaklaşımda hata sınırı]
$f$'nin birinci ve ikinci kısmi türevleri, $(x_0,y_0)$ merkezli bir $B$ dikdörtgeninde sürekli olsun ve $M$,
$B$ üzerinde $|f_{xx}|$, $|f_{xy}|$, $|f_{yy}|$ için ortak bir \textbf{üst sınır} olsun. Bu durumda $B$'nin her
noktasında
\[
\boxed{\,|E(x,y)|=|f(x,y)-L(x,y)|\le\frac{M}{2}\big(|x-x_0|+|y-y_0|\big)^2\,}
\]
\end{teorem}

\textbf{Neden böyle?} (Kaynaklara göre kısa gerekçe.) $P(x_0,y_0)$'dan $(x,y)$'ye giden doğru parçası üzerinde
iki değişkenli Taylor formülü uygulanırsa, parça üzerindeki bir ara noktada hesaplanmak üzere
\[
E(x,y)=\tfrac12\Big[f_{xx}\,(\Delta x)^2+2f_{xy}\,\Delta x\,\Delta y+f_{yy}\,(\Delta y)^2\Big]
\]
bulunur. Mutlak değer alıp her ikinci türevi $M$ ile sınırlarsak
$|E|\le\tfrac12M\big(|\Delta x|^2+2|\Delta x||\Delta y|+|\Delta y|^2\big)=\tfrac12M(|\Delta x|+|\Delta y|)^2$.
Bu da hata formülünün tam olarak tek değişkendeki $|E|\le\frac M2(x-x_0)^2$ formülünün karşılığı olduğunu gösterir.

\begin{bilgi}[Tahtadaki ifade hakkında]
Tahtada ``en büyük değeri $m$'' (\foto{IMG-20261022-WA0060}) / ``$M$'' (\foto{IMG-20261108-WA0093}) denmiştir. Kesin
söylemek gerekirse $M$'nin, ikinci türevlerin mutlak değerlerinin $B$ üzerindeki \emph{herhangi bir üst sınırı} olması
yeterlidir; en küçük üst sınır seçilirse en iyi (en sıkı) tahmin elde edilir. Ayrıca tahtada $B$ açık dikdörtgen
($<h$) olarak tanımlanmış, örneklerde ise kapalı ($\le$) dikdörtgen kullanılmıştır; ikinci türevler sürekli olduğundan
bu fark sonucu etkilemez.
\end{bilgi}

\begin{teorem}[Hata sınırı için adım adım yöntem]
\begin{enumerate}
  \item $f(x_0,y_0)$, $f_x(x_0,y_0)$, $f_y(x_0,y_0)$'ı hesapla ve $L(x,y)$'yi yaz.
  \item $f_{xx}$, $f_{xy}$, $f_{yy}$'yi \textbf{genel} olarak hesapla.
  \item Bu üç fonksiyonun mutlak değerlerini $B$'nin \textbf{tamamında} üstten sınırla ($|\sin u|\le|u|$, $|\cos u|\le1$,
        $|x-x_0|\le h$ gibi eşitsizliklerle); en büyüğünü $M$ al.
  \item $|E|\le\frac M2\big(|x-x_0|+|y-y_0|\big)^2$ ifadesinde $|x-x_0|$ ve $|y-y_0|$ yerine $B$'deki en büyük değerlerini koy.
\end{enumerate}
\end{teorem}

\begin{soru}[Örnek 18 \foto{IMG-20261022-WA0060}, \foto{IMG-20261108-WA0093}]
$f(x,y)=x^2+3xy+2y^2+4$ fonksiyonunun $(2,1)$ noktasındaki lineerleştirmesini bulunuz.
$B=\{(x,y):|x-2|\le0{,}01,\ |y-1|\le0{,}01\}$ dikdörtgensel bölgesi üzerindeki hatanın bir üst sınırını hesaplayınız.
\end{soru}
\begin{cevap}
\textbf{Lineerleştirme.} $f(2,1)=4+6+2+4=16$.\\
$f_x(x,y)=2x+3y\Rightarrow f_x(2,1)=7$;\quad $f_y(x,y)=3x+4y\Rightarrow f_y(2,1)=10$.
\[ L(x,y)=16+7(x-2)+10(y-1)=7x+10y-8 . \]
\textbf{Hata sınırı.} $f_{xx}=2$, $f_{xy}=3$, $f_{yy}=4$ (sabit) $\Rightarrow M=4$.
\[
|E(x,y)|\le\frac12\cdot4\,\big(|x-2|+|y-1|\big)^2\le2\,(0{,}01+0{,}01)^2=2\cdot0{,}0004=0{,}0008 .
\]
\textbf{Yorum.} $B$ üzerinde $f$ yerine $L$ kullanırsak yaptığımız hata en fazla $0{,}0008$'dir. (Bu $f$ ikinci dereceden
bir polinom olduğu için hatanın tam ifadesi $E=(x-2)^2+3(x-2)(y-1)+2(y-1)^2$'dir ve $B$'de en fazla
$0{,}0001+0{,}0003+0{,}0002=0{,}0006$ olur; $0{,}0008$ geçerli bir üst sınırdır.)
\end{cevap}

\begin{soru}[Örnek 19 \foto{IMG-20261105-WA0038}, \foto{IMG-20261105-WA0039}, \foto{IMG-20261108-WA0085}, \foto{IMG-20261108-WA0086}]
$f(x,y)=x^2-3xy+y^2$ fonksiyonunun $P(1,2)$ noktasındaki $L(x,y)$ lineerleştirmesini ve
$B=\{(x,y):|x-1|\le0{,}1,\ |y-2|\le0{,}1\}$ bölgesindeki hatanın bir üst sınırını bulunuz.
\end{soru}
\begin{cevap}
$f(1,2)=1-6+4=-1$;\quad $f_x=2x-3y\Rightarrow f_x(1,2)=-4$;\quad $f_y=-3x+2y\Rightarrow f_y(1,2)=1$.
\[ L(x,y)=-1-4(x-1)+1\cdot(y-2)=-4x+y+1 . \]
İkinci türevler: $f_{xx}=2$, $f_{xy}=-3$, $f_{yy}=2$. $M$ mutlak değerlerin en büyüğüdür:
$M=\max\{2,\,|-3|,\,2\}=3$.
\[ |E(x,y)|\le\frac12\cdot3\,(0{,}1+0{,}1)^2=1{,}5\cdot0{,}04=0{,}06 . \]
\end{cevap}

\begin{soru}[Örnek 20 \foto{IMG-20261105-WA0040}, \foto{IMG-20261108-WA0086}]
$f(x,y)=1+y+x\cos y$ fonksiyonunun $P(0,0)$ noktasındaki lineerleştirmesini ve
$B=\{(x,y):|x|\le0{,}2,\ |y|\le0{,}2\}$ bölgesindeki hatanın bir üst sınırını bulunuz.
\end{soru}
\begin{cevap}
$f(0,0)=1$;\quad $f_x=\cos y\Rightarrow f_x(0,0)=1$;\quad $f_y=1-x\sin y\Rightarrow f_y(0,0)=1$.
\[ L(x,y)=1+x+y . \]
İkinci türevler: $f_{xx}=0$, $f_{xy}=-\sin y$, $f_{yy}=-x\cos y$. Bunları $B$'nin tamamında sınırlayalım:
\[
|f_{xy}|=|\sin y|\le|y|\le0{,}2,\qquad |f_{yy}|=|x|\,|\cos y|\le0{,}2\cdot1=0{,}2\quad\Rightarrow\quad M=0{,}2 .
\]
\[ |E(x,y)|\le\frac12\cdot0{,}2\,(0{,}2+0{,}2)^2=0{,}1\cdot0{,}16=0{,}016 . \]
\textbf{Dikkat:} İkinci türevler $P(0,0)$'da hesaplansaydı hepsi $0$ çıkardı ve yanlışlıkla $M=0$ bulunurdu. $M$,
$B$'nin tamamı üzerindeki üst sınırdır.
\end{cevap}

\begin{soru}[\ekal Örnek 21 (üstel fonksiyonda hata sınırı)]
$f(x,y)=e^x\cos y$'nin $(0,0)$'daki lineerleştirmesini ve $B:\ |x|\le0{,}1,\ |y|\le0{,}1$ üzerindeki hata sınırını bulunuz.
\end{soru}
\begin{cevap}
$f(0,0)=1$, $f_x=e^x\cos y\Rightarrow1$, $f_y=-e^x\sin y\Rightarrow0$; $L(x,y)=1+x$.\\
$f_{xx}=e^x\cos y$, $f_{xy}=-e^x\sin y$, $f_{yy}=-e^x\cos y$. $B$'de $|\cos y|\le1$, $|\sin y|\le1$ ve $e^x\le e^{0{,}1}<1{,}11$
olduğundan $M=1{,}11$ alınabilir:
\[ |E|\le\frac{1{,}11}{2}\,(0{,}1+0{,}1)^2=0{,}0222 . \]
(Daha sıkı: $|f_{xy}|\le e^{0{,}1}\sin0{,}1\approx0{,}11$ ama $|f_{xx}|$ yine $\approx1{,}105$'e ulaşır; $M$ en büyük olana göre seçilir.)
\end{cevap}

%-------------------------------------------------------------------------------------------
\subsection{Tam diferansiyel ve yüksek mertebeden diferansiyeller}

\subsubsection{Birinci mertebe tam (toplam) diferansiyel}

\begin{tahta}[\foto{IMG-20261108-WA0065}]
\textbf{Tam Diferansiyel:}
\[
\Delta z=f(x,y)-f(x_0,y_0)=f_x(x_0,y_0)\Delta x+f_y(x_0,y_0)\Delta y+\varepsilon_1\Delta x+\varepsilon_2\Delta y
\]
$f_x\Delta x$ ve $f_y\Delta y$: $f$ fonksiyonunun $x$ ve $y$'ye göre \textbf{kısmi diferansiyelleri}.\\
$dz=f_x\Delta x+f_y\Delta y$: \textbf{tam (toplam) diferansiyel}.\\
$f(x,y)=x$ olsaydı $\Delta x=dx$; $f(x,y)=y$ olsaydı $\Delta y=dy$ $\Rightarrow\ dz=f_x\,dx+f_y\,dy$.\\
Üç değişkende: $w=f(x,y,z)\Rightarrow dw=f_x\,dx+f_y\,dy+f_z\,dz$.
\end{tahta}

\begin{tanim}[Tam diferansiyel]
$z=f(x,y)$ diferansiyellenebilir olsun. $dx=\Delta x$ ve $dy=\Delta y$ bağımsız değişkenlerin (keyfi) artımları olmak
üzere
\[ \boxed{\,dz=f_x(x,y)\,dx+f_y(x,y)\,dy=\frac{\partial z}{\partial x}dx+\frac{\partial z}{\partial y}dy\,} \]
ifadesine $z$'nin \textbf{tam (toplam) diferansiyeli} denir. $f_x\,dx$ ve $f_y\,dy$ terimleri \textbf{kısmi
diferansiyellerdir}.
\end{tanim}

Tahtadaki ``$f(x,y)=x$ olsaydı'' açıklamasının anlamı şudur: $f(x,y)=x$ fonksiyonu için $f_x=1$, $f_y=0$ olduğundan
$dx=1\cdot\Delta x+0\cdot\Delta y=\Delta x$ olur. Yani bağımsız değişkenler için diferansiyel ile artım aynıdır; bu yüzden
$\Delta x$ yerine $dx$ yazabiliriz.

\textbf{$dz$ ile $\Delta z$ arasındaki fark.} $\Delta z$, $f$'nin \emph{gerçek} değişimidir (yüzey üzerinde);
$dz$ ise teğet düzlem üzerindeki değişimdir. Artım formülünden
\[ \Delta z-dz=\varepsilon_1\,dx+\varepsilon_2\,dy\ \longrightarrow\ 0 \]
ve bu fark $dx,dy$'den çok daha hızlı küçülür. Bu yüzden küçük artımlar için $\Delta z\approx dz$. Tahtadaki çizim, tek
değişkenli durumda bunu gösteriyordu:

\begin{center}
\begin{tikzpicture}[scale=1.25]
\draw[eks] (-0.2,0) -- (4,0) node[below]{$x$};
\draw[eks] (0,-0.2) -- (0,3) node[left]{$y$};
\draw[very thick,ana,domain=0.2:3.4,samples=25] plot (\x,{0.25*\x*\x+0.5}) node[above]{$y=f(x)$};
\draw[thick,vurgu,domain=0.2:3.4] plot (\x,{0.86+0.6*(\x-1.2)});
\node[vurgu,below right] at (0.2,0.26) {teğet};
\fill (1.2,0.86) circle (1.5pt);
\draw[dashed,gri] (1.2,0.86) -- (1.2,0) node[below,black]{$x_0$};
\draw[dashed,gri] (2.8,2.46) -- (2.8,0) node[below,black]{$x_0+dx$};
\draw[dashed,gri] (1.2,0.86) -- (2.8,0.86);
\draw[-{Stealth[length=2mm]},thick] (1.2,0.55) -- node[below]{$dx=\Delta x$} (2.8,0.55);
\draw[-{Stealth[length=2mm]},very thick,vurgu] (2.95,0.86) -- node[right]{$dy$} (2.95,1.82);
\draw[-{Stealth[length=2mm]},very thick,ana] (3.5,0.86) -- node[right]{$\Delta y$} (3.5,2.46);
\draw[dashed,gri] (2.8,0.86) -- (3.55,0.86);
\draw[dashed,gri] (2.8,2.46) -- (3.55,2.46);
\draw[dashed,gri] (2.8,1.82) -- (3.0,1.82);
\end{tikzpicture}
\end{center}
\begin{center}
{\small Şekil 4. \foto{IMG-20261108-WA0065} üzerindeki çizimin yeniden çizimi: $dy$ teğet doğru boyunca artış,
$\Delta y$ eğri boyunca gerçek artış; $dx$ küçüldükçe $dy\approx\Delta y$.}
\end{center}

\begin{soru}[Örnek 22 \foto{IMG-20261105-WA0034}, \foto{IMG-20261108-WA0082} (Problem 1a, birinci mertebe)]
$f(x,y)=x^3+5xy^2+3xy+4x-3y+5$ fonksiyonunun birinci mertebeden tam diferansiyelini bulunuz.
\end{soru}
\begin{cevap}
$f_x=3x^2+5y^2+3y+4$,\quad $f_y=10xy+3x-3$. Dolayısıyla
\[ df=(3x^2+5y^2+3y+4)\,dx+(10xy+3x-3)\,dy . \]
(İkinci mertebe kısmı aşağıda, Örnek 26'da.)
\end{cevap}

\begin{soru}[Örnek 23 \foto{IMG-20261105-WA0034}, \foto{IMG-20261108-WA0082} (Problem 1b, birinci mertebe)]
$w=f(x,y,z)=3x^2yz^2-3xy+2yz^2+xyz$ fonksiyonunun birinci mertebeden tam diferansiyelini bulunuz.
\end{soru}
\begin{cevap}
Her değişkene göre, diğer ikisi sabit tutularak:
\[
w_x=6xyz^2-3y+yz,\qquad w_y=3x^2z^2-3x+2z^2+xz,\qquad w_z=6x^2yz+4yz+xy .
\]
\[
dw=(6xyz^2-3y+yz)\,dx+(3x^2z^2-3x+2z^2+xz)\,dy+(6x^2yz+4yz+xy)\,dz .
\]
\end{cevap}

\begin{soru}[\ekal Örnek 24 ($\Delta z$ ile $dz$'nin karşılaştırılması)]
$z=x^2+3xy-y^2$ için $(x,y)$, $(2,3)$'ten $(2{,}05;\,2{,}96)$'ya değiştiğinde $dz$ ve $\Delta z$'yi hesaplayıp
karşılaştırınız.
\end{soru}
\begin{cevap}
$dz=(2x+3y)\,dx+(3x-2y)\,dy$. $x=2$, $y=3$, $dx=0{,}05$, $dy=-0{,}04$:
\[ dz=(4+9)(0{,}05)+(6-6)(-0{,}04)=0{,}65 . \]
Gerçek değişim: $f(2,3)=4+18-9=13$, $f(2{,}05;\,2{,}96)=4{,}2025+18{,}204-8{,}7616=13{,}6449$ ve
$\Delta z=0{,}6449$. Fark yalnızca $0{,}0051$; $dz$'yi hesaplamak çok daha kolaydır.
\end{cevap}

\subsubsection{İkinci ve yüksek mertebeden diferansiyeller}

\begin{tahta}[\foto{IMG-20261108-WA0068}]
\textbf{Yüksek Mertebeden Diferansiyeller:} $z=f(x,y)$, $dz=z_x\,dx+z_y\,dy$. $dx$ ve $dy$ sabit artımlar kabul edilip
bir kez daha diferansiyel alınırsa:
\[ d^2z=d(dz)=z_{xx}\,dx^2+2z_{xy}\,dx\,dy+z_{yy}\,dy^2 \]
Sembolik yazılış:
\[
d^2z=\left(\frac{\partial}{\partial x}dx+\frac{\partial}{\partial y}dy\right)^{(2)}z,\qquad
d^nz=\left(\frac{\partial}{\partial x}dx+\frac{\partial}{\partial y}dy\right)^{(n)}z
\]
Üç değişkende: $d^2w=w_{xx}dx^2+w_{yy}dy^2+w_{zz}dz^2+2w_{xy}dx\,dy+2w_{xz}dx\,dz+2w_{yz}dy\,dz$.
\end{tahta}

\textbf{Formül nereden geliyor?} $dx$, $dy$ sabit olduğundan $d(dz)=d(z_x)\,dx+d(z_y)\,dy$ ve
$d(z_x)=z_{xx}dx+z_{xy}dy$, $d(z_y)=z_{yx}dx+z_{yy}dy$. Yerine koyarsak
\[ d^2z=z_{xx}\,dx^2+(z_{xy}+z_{yx})\,dx\,dy+z_{yy}\,dy^2 . \]
Schwarz teoremi gereği (ikinci türevler sürekli ise) $z_{xy}=z_{yx}$ olduğundan orta terim $2z_{xy}\,dx\,dy$ olur. Tahtada
bu varsayım açıkça yazılmamıştı; formül, ikinci kısmi türevlerin sürekli olduğu durumda geçerlidir.

\begin{teorem}[Yüksek mertebeden diferansiyeller (binom açılımı gibi)]
Gerekli kısmi türevler sürekli ise, sembolik kuvvet binom formülüyle açılır; $dx^2$, $(dx)^2$ demektir:
\begin{align*}
d^2z&=z_{xx}\,dx^2+2z_{xy}\,dx\,dy+z_{yy}\,dy^2,\\
d^3z&=z_{xxx}\,dx^3+3z_{xxy}\,dx^2dy+3z_{xyy}\,dx\,dy^2+z_{yyy}\,dy^3,\\
d^nz&=\sum_{k=0}^{n}\binom nk\frac{\partial^nz}{\partial x^{n-k}\,\partial y^{k}}\,dx^{n-k}dy^{k}.
\end{align*}
\end{teorem}

\begin{soru}[Örnek 25 \foto{IMG-20261108-WA0069}]
$z=x^2y$ fonksiyonunun ikinci mertebeden diferansiyelini bulunuz.
\end{soru}
\begin{cevap}
$z_x=2xy$, $z_y=x^2$;\quad $z_{xx}=2y$, $z_{xy}=2x$, $z_{yy}=0$.
\[ d^2z=2y\,dx^2+2\cdot2x\,dx\,dy+0\cdot dy^2=2y\,dx^2+4x\,dx\,dy . \]
\end{cevap}

\begin{soru}[Örnek 26 \foto{IMG-20261105-WA0034}, \foto{IMG-20261108-WA0082} (Problem 1a, ikinci mertebe)]
$f(x,y)=x^3+5xy^2+3xy+4x-3y+5$ için $d^2f$'yi bulunuz.
\end{soru}
\begin{cevap}
Örnek 22'den $f_x=3x^2+5y^2+3y+4$, $f_y=10xy+3x-3$. Buradan
$f_{xx}=6x$,\quad $f_{xy}=10y+3$ $(=f_{yx}=(10xy+3x-3)_x$ $\checkmark$),\quad $f_{yy}=10x$.
\[ d^2f=6x\,dx^2+2(10y+3)\,dx\,dy+10x\,dy^2 . \]
\end{cevap}

\begin{soru}[Örnek 27 \foto{IMG-20261105-WA0035}, \foto{IMG-20261108-WA0083} (Problem 1b, ikinci mertebe)]
$w=3x^2yz^2-3xy+2yz^2+xyz$ için $d^2w$'yi bulunuz.
\end{soru}
\begin{cevap}
Örnek 23'teki birinci türevlerden:
\begin{align*}
w_{xx}&=6yz^2, & w_{yy}&=0, & w_{zz}&=6x^2y+4y,\\
w_{xy}&=6xz^2-3+z, & w_{xz}&=12xyz+y, & w_{yz}&=6x^2z+4z+x .
\end{align*}
Üç değişkenli formüle koyarsak ($w_{yy}=0$ olduğundan $dy^2$ terimi yoktur):
\begin{align*}
d^2w={}&6yz^2\,dx^2+(6x^2y+4y)\,dz^2+2(6xz^2+z-3)\,dx\,dy\\
&+2(12xyz+y)\,dx\,dz+2(6x^2z+x+4z)\,dy\,dz .
\end{align*}
\end{cevap}

\begin{soru}[\ekal Örnek 28]
$z=e^{xy}$ için $d^2z$'yi bulunuz.
\end{soru}
\begin{cevap}
$z_x=y\,e^{xy}$, $z_y=x\,e^{xy}$; $z_{xx}=y^2e^{xy}$, $z_{yy}=x^2e^{xy}$, $z_{xy}=e^{xy}+xy\,e^{xy}=(1+xy)e^{xy}$.
\[ d^2z=e^{xy}\big[y^2\,dx^2+2(1+xy)\,dx\,dy+x^2\,dy^2\big]. \]
\end{cevap}

%-------------------------------------------------------------------------------------------
\subsection{Diferansiyelle yaklaşık hesap: mutlak, bağıl ve yüzde değişim}

Tam diferansiyelin en önemli uygulaması, bağımsız değişkenler küçük miktarlarda değiştiğinde (ya da küçük ölçüm
hatalarıyla ölçüldüğünde) fonksiyondaki değişimi kolayca tahmin etmektir:
\[ \Delta z\approx dz=f_x\,dx+f_y\,dy,\qquad f(x_0+dx,\,y_0+dy)\approx f(x_0,y_0)+dz . \]

\begin{tanim}[Mutlak, bağıl ve yüzde değişim]
\begin{itemize}
  \item \textbf{Mutlak (yaklaşık) değişim:} $\Delta z\approx dz$.
  \item \textbf{Bağıl değişim:} $\dfrac{\Delta z}{z}\approx\dfrac{dz}{z}$ (başlangıç değerine oranla değişim).
  \item \textbf{Yüzde değişim:} $\dfrac{dz}{z}\cdot100$ \ (\%).
\end{itemize}
Ölçüm hatası problemlerinde $dx$, $dy$ ölçümdeki en büyük hatalardır; en kötü durumu bulmak için
$|dz|\le|f_x||dx|+|f_y||dy|$ kullanılır (\textbf{en büyük mutlak hata}); $|dz|/|z|$ en büyük bağıl hatadır.
\end{tanim}

\begin{teorem}[Logaritmik diferansiyel kısayolu]
$z=C\,x^a y^b$ biçimindeki (çarpım/kuvvet) fonksiyonlarda bağıl değişim doğrudan
\[ \frac{dz}{z}=a\,\frac{dx}{x}+b\,\frac{dy}{y} \]
ile bulunur ($\ln z=\ln C+a\ln x+b\ln y$'nin diferansiyeli). Örneğin $V=\pi r^2h$ için
$\dfrac{dV}{V}=2\dfrac{dr}{r}+\dfrac{dh}{h}$.
\end{teorem}

\begin{teorem}[Adım adım yöntem]
\begin{enumerate}
  \item Büyüklüğü bağımsız değişkenler cinsinden yaz (ör. $V=\pi r^2h$).
  \item Tam diferansiyeli yaz: $dV=V_r\,dr+V_h\,dh$.
  \item Başlangıç değerlerini ve artımları belirle; \textbf{azalma ise artım negatiftir}.
  \item Sayısal değerleri yerleştir; istenirse yeni değeri $V+dV$, bağıl/yüzde değişimi $dV/V$ ile bul.
\end{enumerate}
\end{teorem}

\begin{soru}[Örnek 29 \foto{IMG-20261108-WA0065}, \foto{IMG-20261108-WA0066}, \foto{IMG-20261108-WA0067}]
Yarıçapı $r=3$ cm, yüksekliği $h=10$ cm olan bir dik dairesel silindirin yarıçapı $0{,}01$ cm arttırılıp yüksekliği
$0{,}1$ cm azaltılırsa hacmindeki değişim nedir?
\end{soru}
\begin{cevap}
$V=\pi r^2h\Rightarrow dV=V_r\,dr+V_h\,dh=2\pi rh\,dr+\pi r^2\,dh$.\quad $dr=0{,}01$, $dh=-0{,}1$ (azalma).
\[
dV=2\pi\cdot3\cdot10\,(0{,}01)+\pi\cdot9\,(-0{,}1)=0{,}6\pi-0{,}9\pi=-0{,}3\pi\ \text{cm}^3\approx-0{,}94\ \text{cm}^3 .
\]
Hacim yaklaşık $0{,}3\pi$ cm$^3$ azalır. Başlangıç hacmi $V=\pi\cdot9\cdot10=90\pi$, dolayısıyla
\[ \text{yeni hacim}\approx90\pi-0{,}3\pi=89{,}7\pi\ \text{cm}^3\approx281{,}8\ \text{cm}^3 . \]
(Kontrol: gerçek yeni hacim $\pi\cdot3{,}01^2\cdot9{,}9\approx89{,}695\pi$ cm$^3$.)
\end{cevap}

\begin{dikkat}[Tahtada düzeltilen hata]
\foto{IMG-20261108-WA0066} ve \foto{IMG-20261108-WA0067}'de son satır ``yeni hacim $\approx90\pi-0{,}3\pi=89{,}7$ cm$^3$''
olarak yazılmıştı; $\pi$ düşmüştür. Doğrusu $89{,}7\pi$ cm$^3\approx281{,}8$ cm$^3$'tür.
\end{dikkat}

\begin{soru}[Örnek 30 \foto{IMG-20261105-WA0035}, \foto{IMG-20261108-WA0083}, \foto{IMG-20261108-WA0084} (Problem 2)]
Bir dik dairesel silindirin yarıçapı 10 cm'den 9,9 cm'ye azalıyor, yüksekliği 30 cm'den 30,3 cm'ye çıkıyor. Hacimdeki
değişimi diferansiyel ile yaklaşık olarak hesaplayınız.
\end{soru}
\begin{cevap}
$dV=2\pi rh\,dr+\pi r^2\,dh$;\quad $r=10$, $h=30$, $dr=9{,}9-10=-0{,}1$, $dh=30{,}3-30=0{,}3$.
\[
dV=2\pi\cdot10\cdot30\cdot(-0{,}1)+\pi\cdot100\cdot0{,}3=-60\pi+30\pi=-30\pi\ \text{cm}^3\ (\approx-94{,}2\ \text{cm}^3).
\]
$V=\pi\cdot100\cdot30=3000\pi$ olduğundan yeni hacim $\approx3000\pi-30\pi=2970\pi$ cm$^3$.
(Kontrol: gerçek değer $\pi\cdot9{,}9^2\cdot30{,}3\approx2969{,}70\pi$ cm$^3$; yaklaşım çok iyi.)
\end{cevap}

\begin{soru}[Örnek 31 \foto{IMG-20261105-WA0036} (Problem 3)]
Yarıçapı $r=1$, yüksekliği $h=5$ olan bir silindirin yarıçapı $0{,}03$ arttırılıp, yüksekliği $0{,}1$ azaltılırsa
hacimdeki yaklaşık değişim ve yüzde değişim nedir?
\end{soru}
\begin{cevap}
$dV=2\pi rh\,dr+\pi r^2\,dh$; $dr=0{,}03$, $dh=-0{,}1$:
\[ dV=2\pi\cdot1\cdot5\cdot0{,}03+\pi\cdot1\cdot(-0{,}1)=0{,}3\pi-0{,}1\pi=0{,}2\pi . \]
$V=\pi\cdot1\cdot5=5\pi$;\quad $\dfrac{dV}{V}=\dfrac{0{,}2\pi}{5\pi}=0{,}04$. Hacim yaklaşık \textbf{\%4 artar}.\\
Logaritmik kısayolla: $\dfrac{dV}{V}=2\dfrac{dr}{r}+\dfrac{dh}{h}=2\cdot\dfrac{0{,}03}{1}+\dfrac{-0{,}1}{5}=0{,}06-0{,}02=0{,}04$ $\checkmark$.
(Gerçek değişim $\pi(1{,}03^2\cdot4{,}9-5)\approx0{,}198\pi$, yani yaklaşık \%3,97.)
\end{cevap}

\begin{soru}[\ekal Örnek 32 (yaklaşık değer)]
$(0{,}98)^2\sqrt{4{,}03}$ sayısını diferansiyel yardımıyla yaklaşık hesaplayınız.
\end{soru}
\begin{cevap}
$f(x,y)=x^2\sqrt y$, $(x_0,y_0)=(1,4)$, $dx=-0{,}02$, $dy=0{,}03$. $f(1,4)=2$;
$f_x=2x\sqrt y\Rightarrow4$; $f_y=\dfrac{x^2}{2\sqrt y}\Rightarrow\dfrac14$.
\[ df=4(-0{,}02)+\tfrac14(0{,}03)=-0{,}08+0{,}0075=-0{,}0725,\qquad f\approx2-0{,}0725=1{,}9275 . \]
(Hesap makinesi: $1{,}92799\ldots$)
\end{cevap}

\begin{soru}[\ekal Örnek 33 (ölçüm hatası: koni)]
Bir dik koninin taban yarıçapı ve yüksekliği 10 cm ve 25 cm olarak ölçülmüş; her ölçümde en fazla $0{,}1$ cm hata
olabilir. Hesaplanan hacimdeki en büyük hatayı ve en büyük bağıl hatayı tahmin ediniz.
\end{soru}
\begin{cevap}
$V=\dfrac13\pi r^2h$, $dV=\dfrac{2\pi rh}{3}dr+\dfrac{\pi r^2}{3}dh$. En kötü durumda iki hata aynı işaretlidir:
$|dr|,|dh|\le0{,}1$,
\[
|dV|\le\frac{2\pi\cdot10\cdot25}{3}(0{,}1)+\frac{\pi\cdot100}{3}(0{,}1)=\frac{50\pi}{3}+\frac{10\pi}{3}=20\pi\approx62{,}8\ \text{cm}^3 .
\]
$V=\dfrac{2500\pi}{3}$, bağıl hata $\dfrac{20\pi}{2500\pi/3}=0{,}024$, yani en fazla yaklaşık \%2,4.
Logaritmik kısayolla da: $\dfrac{|dV|}{V}\le2\dfrac{0{,}1}{10}+\dfrac{0{,}1}{25}=0{,}02+0{,}004=0{,}024$ $\checkmark$.
\end{cevap}

\begin{soru}[\ekal Örnek 34 (yüzde hata: sarkaç)]
Basit sarkacın periyodu $T=2\pi\sqrt{L/g}$. $L$ en fazla \%1, $g$ en fazla \%0,5 hatayla biliniyorsa $T$'deki en büyük
yüzde hata nedir?
\end{soru}
\begin{cevap}
$\ln T=\ln2\pi+\tfrac12\ln L-\tfrac12\ln g\Rightarrow \dfrac{dT}{T}=\dfrac12\dfrac{dL}{L}-\dfrac12\dfrac{dg}{g}$.
En kötü durumda terimler aynı işaretli olur:
$\left|\dfrac{dT}{T}\right|\le\dfrac12(0{,}01)+\dfrac12(0{,}005)=0{,}0075$, yani en fazla \%0,75.
\end{cevap}

%-------------------------------------------------------------------------------------------
\subsection{Üç değişkenli fonksiyonlara genelleme}

\begin{bilgi}[Kısmen tahtada]
Tahtada üç değişkenli fonksiyonlar için tam diferansiyel $dw$ ve $d^2w$ formülleri ile Problem 1b işlendi
(\foto{IMG-20261108-WA0065}, \foto{IMG-20261108-WA0068}). Kısmi türev tanımı, diferansiyellenebilirlik, lineerleştirme ve
hata sınırının üç değişkenli hâlleri kaynaklara göre eklenmiştir.
\end{bilgi}

$w=f(x,y,z)$ fonksiyonunun kısmi türevleri, diğer \emph{iki} değişken sabit tutularak tanımlanır:
\begin{align*}
f_x&=\lim_{h\to0}\frac{f(x+h,y,z)-f(x,y,z)}{h},\qquad
f_y=\lim_{h\to0}\frac{f(x,y+h,z)-f(x,y,z)}{h},\\
f_z&=\lim_{h\to0}\frac{f(x,y,z+h)-f(x,y,z)}{h}.
\end{align*}
Bu kez $3^2=9$ ikinci mertebe kısmi türev vardır; bunlar sürekli ise Schwarz teoremi gereği
$f_{xy}=f_{yx}$, $f_{xz}=f_{zx}$, $f_{yz}=f_{zy}$ olduğundan yalnızca $6$ tanesi farklıdır.

\begin{teorem}[Üç değişkende formüller]
\begin{description}
  \item[Artım ve diferansiyellenebilirlik:]
  $\Delta w=f_x\Delta x+f_y\Delta y+f_z\Delta z+\varepsilon_1\Delta x+\varepsilon_2\Delta y+\varepsilon_3\Delta z$,
  $\varepsilon_i\to0$. Kısmi türevler sürekli ise $f$ diferansiyellenebilirdir.
  \item[Tam diferansiyel:] $dw=f_x\,dx+f_y\,dy+f_z\,dz$.
  \item[İkinci diferansiyel:] $d^2w=w_{xx}dx^2+w_{yy}dy^2+w_{zz}dz^2+2w_{xy}dx\,dy+2w_{xz}dx\,dz+2w_{yz}dy\,dz$.
  \item[Lineerleştirme:]
  $L(x,y,z)=f(P_0)+f_x(P_0)(x-x_0)+f_y(P_0)(y-y_0)+f_z(P_0)(z-z_0)$, \ $P_0=(x_0,y_0,z_0)$.
  \item[Hata sınırı:] İkinci türevler $P_0$ merkezli kutu biçimli $B$ bölgesinde sürekli ve hepsinin mutlak değeri
  $\le M$ ise
  \[ |E(x,y,z)|\le\frac M2\big(|x-x_0|+|y-y_0|+|z-z_0|\big)^2 . \]
\end{description}
\end{teorem}

Not: Üç değişkenli bir fonksiyonun grafiği dört boyutlu olduğundan ``teğet düzlem'' yerine burada yalnızca lineer yaklaşım
konuşulur. $F(x,y,z)=0$ seviye yüzeyinin teğet düzlemi gradyan ile bulunur; bkz. Konu 6.

\begin{soru}[\ekal Örnek 35 (üç değişkende yaklaşık hesap)]
$\sqrt{(3{,}02)^2+(1{,}97)^2+(5{,}99)^2}$ sayısını lineer yaklaşımla hesaplayınız.
\end{soru}
\begin{cevap}
$f(x,y,z)=\sqrt{x^2+y^2+z^2}$, $P_0=(3,2,6)$: $f(P_0)=\sqrt{49}=7$.
$f_x=\dfrac xf$, $f_y=\dfrac yf$, $f_z=\dfrac zf$ $\Rightarrow$ $P_0$'da $\dfrac37,\ \dfrac27,\ \dfrac67$.
$dx=0{,}02$, $dy=-0{,}03$, $dz=-0{,}01$:
\[
df=\frac{3(0{,}02)+2(-0{,}03)+6(-0{,}01)}{7}=\frac{0{,}06-0{,}06-0{,}06}{7}\approx-0{,}00857,\qquad f\approx6{,}9914 .
\]
(Gerçek değer $\approx6{,}99152$.)
\end{cevap}

\begin{soru}[\ekal Örnek 36 (üç değişkende lineerleştirme ve hata sınırı)]
$f(x,y,z)=xz-3yz+2$ fonksiyonunun $P_0(1,1,2)$'deki lineerleştirmesini ve
$B:\ |x-1|\le0{,}01,\ |y-1|\le0{,}01,\ |z-2|\le0{,}02$ kutusundaki hata sınırını bulunuz.
\end{soru}
\begin{cevap}
$f(P_0)=2-6+2=-2$;\ $f_x=z\Rightarrow2$;\ $f_y=-3z\Rightarrow-6$;\ $f_z=x-3y\Rightarrow-2$.
\[ L=-2+2(x-1)-6(y-1)-2(z-2)=2x-6y-2z+6 . \]
İkinci türevler: $f_{xx}=f_{yy}=f_{zz}=0$, $f_{xy}=0$, $f_{xz}=1$, $f_{yz}=-3$ $\Rightarrow M=3$.
\[ |E|\le\frac32\,(0{,}01+0{,}01+0{,}02)^2=1{,}5\cdot0{,}0016=0{,}0024 . \]
\end{cevap}

\begin{soru}[\ekal Örnek 37 (üç değişkende tam diferansiyel)]
$w=x^2y+y^2z+\ln(xz)$ $(x,z>0)$ için $dw$'yi ve $(1,1,1)$ noktasında $dx=0{,}1$, $dy=-0{,}1$, $dz=0{,}2$ için değerini bulunuz.
\end{soru}
\begin{cevap}
$w_x=2xy+\dfrac1x$, $w_y=x^2+2yz$, $w_z=y^2+\dfrac1z$. Böylece
$dw=\left(2xy+\dfrac1x\right)dx+(x^2+2yz)\,dy+\left(y^2+\dfrac1z\right)dz$.
$(1,1,1)$'de katsayılar $3,\ 3,\ 2$: $dw=3(0{,}1)+3(-0{,}1)+2(0{,}2)=0{,}4$.
\end{cevap}

%-------------------------------------------------------------------------------------------
\subsection{Kaynak kitaplardan ek örnekler}

Aşağıdaki örnekler, konunun farklı parçalarını bir arada kullanan, kaynak kitaplardaki tipik problemlere benzer
örneklerdir.

\begin{soru}[\ekal Ek örnek A (teğet düzlem + yaklaşık değer)]
$f(x,y)=\sqrt{20-x^2-7y^2}$ yüzeyinin $(2,1)$ noktasındaki teğet düzlemini bulunuz ve $f(1{,}95;\,1{,}08)$'i yaklaşık hesaplayınız.
\end{soru}
\begin{cevap}
$f(2,1)=\sqrt{20-4-7}=3$. $f_x=\dfrac{-x}{\sqrt{20-x^2-7y^2}}\Rightarrow-\dfrac23$;\
$f_y=\dfrac{-7y}{\sqrt{20-x^2-7y^2}}\Rightarrow-\dfrac73$.
\[ z=3-\tfrac23(x-2)-\tfrac73(y-1) . \]
$x-2=-0{,}05$, $y-1=0{,}08$: $f(1{,}95;\,1{,}08)\approx3+0{,}0333-0{,}1867=2{,}8467$. (Gerçek değer $\approx2{,}8342$;
$y$'deki artım görece büyük olduğundan hata biraz daha büyüktür.)
\end{cevap}

\begin{soru}[\ekal Ek örnek B (kutu hacmi)]
Boyutları 50 cm, 40 cm, 30 cm olarak ölçülen dikdörtgenler prizması biçimli bir kutunun her boyutunda en fazla
$0{,}1$ cm ölçüm hatası vardır. Hacimdeki en büyük hatayı ve yüzde hatayı tahmin ediniz.
\end{soru}
\begin{cevap}
$V=xyz$, $dV=yz\,dx+xz\,dy+xy\,dz$. $|dx|,|dy|,|dz|\le0{,}1$:
\[ |dV|\le(40\cdot30+50\cdot30+50\cdot40)(0{,}1)=(1200+1500+2000)(0{,}1)=470\ \text{cm}^3 . \]
$V=60000$ cm$^3$; bağıl hata $\le470/60000\approx0{,}0078$, yani yaklaşık \%0,78
(kısayol: $\frac{0{,}1}{50}+\frac{0{,}1}{40}+\frac{0{,}1}{30}\approx0{,}0078$).
\end{cevap}

\begin{soru}[\ekal Ek örnek C (diferansiyel ve teğet düzlem bağlantısı)]
$z=x^2+y^2$ için $(1,2)$ noktasında $dx=0{,}1$, $dy=0{,}05$ alınıyor. $dz$'yi, $\Delta z$'yi bulunuz ve $dz$'nin teğet
düzlemdeki yükseklik değişimi olduğunu gösteriniz.
\end{soru}
\begin{cevap}
$dz=2x\,dx+2y\,dy=2(0{,}1)+4(0{,}05)=0{,}4$. $\Delta z=(1{,}1^2+2{,}05^2)-5=1{,}21+4{,}2025-5=0{,}4125$.
$(1,2,5)$'teki teğet düzlem $z=5+2(x-1)+4(y-2)$; $(1{,}1;\,2{,}05)$'te düzlemin yüksekliği $5+0{,}2+0{,}2=5{,}4$, yani
düzlem üzerindeki değişim tam olarak $0{,}4=dz$. Fark $\Delta z-dz=0{,}0125=dx^2+dy^2$ (Örnek 10(a)'daki
$\varepsilon_1dx+\varepsilon_2dy$ terimi).
\end{cevap}

%-------------------------------------------------------------------------------------------
\subsection{Sık yapılan hatalar}

\begin{dikkat}
\begin{itemize}
  \item \textbf{Diğer değişkeni sabit tutmamak:} $f=x^2y^3$ için $f_x$'i ``$2xy^3+3x^2y^2$'' gibi yazmak. $f_x$'te $y$
        bir sayıdır: $f_x=2xy^3$.
  \item \textbf{Sabit terimin türevi:} $f_y$ hesaplanırken yalnızca $x$ içeren terimlerin ($3x$ gibi) türevi $0$'dır.
  \item \textbf{Karma türevde sıra:} $f_{xy}=\dfrac{\partial^2f}{\partial y\,\partial x}$ (alt indis soldan sağa, kesir
        sağdan sola).
  \item \textbf{Clairaut teoremini yanlış koşulla kullanmak:} Karma türevlerin eşitliği fonksiyonun simetrik olmasına
        değil, karma türevlerin \emph{sürekliliğine} bağlıdır (tahtadaki ``$f(x,y)=f(y,x)$ ise'' notu yanlıştır).
  \item \textbf{``Kısmi türevler var $\Rightarrow$ diferansiyellenebilir'' sanmak:} İki değişkende bu yanlıştır
        (Örnek 11, 12). Yeter koşul kısmi türevlerin \emph{sürekli} olmasıdır.
  \item \textbf{Teğet düzlemde türevleri noktada hesaplamamak:} $z-z_0=f_x(x-x_0)+f_y(y-y_0)$'da $f_x$, $f_y$
        \emph{sayı} olmalıdır; aksi hâlde sonuç bir düzlem olmaz. $z_0=f(x_0,y_0)$'ı unutmamak gerekir.
  \item \textbf{Normal vektörün işareti:} $z=f(x,y)$ için normal $\langle f_x,f_y,-1\rangle$'dir; $+1$ yazmak yanlış düzlem verir.
  \item \textbf{Hata sınırında $M$'yi yalnızca noktada hesaplamak:} $M$, ikinci türevlerin mutlak değerlerinin $B$'nin
        \emph{tamamındaki} üst sınırıdır (Örnek 20'de noktada hesaplansa $M=0$ çıkardı). Negatif türevlerin mutlak değeri
        alınır ($f_{xy}=-3\Rightarrow3$).
  \item \textbf{Azalmayı pozitif almak:} ``yükseklik $0{,}1$ azaltılıyor'' ise $dh=-0{,}1$'dir.
  \item \textbf{Değişim ile yeni değeri karıştırmak:} $dV$ \emph{değişimdir}; yeni değer $V+dV$'dir. $\pi$ gibi çarpanları
        son satırda düşürmemek gerekir (tahtadaki $89{,}7$ cm$^3$ hatası; doğrusu $89{,}7\pi$ cm$^3$).
  \item \textbf{Yüzde değişimde başlangıç değeri:} $dV/V$'de $V$ başlangıç değeridir; sonucu yüzde olarak vermek için
        $100$ ile çarpılır.
  \item \textbf{$d^2z$'de orta terimin katsayısı:} $2z_{xy}\,dx\,dy$; $2$'yi unutmamak. $dx^2$, $(dx)^2$ demektir,
        $d(x^2)$ değil.
  \item \textbf{Ölçüm hatasında en kötü durum:} En büyük hata istenirken terimlerin mutlak değerleri toplanır;
        işaretli toplam, hataların birbirini götürdüğü iyimser bir tahmindir.
\end{itemize}
\end{dikkat}

%-------------------------------------------------------------------------------------------
\subsection{Özet ve formül kutusu}

\begin{ozet}
\textbf{Kısmi türevler}
\[
f_x(a,b)=\lim_{h\to0}\frac{f(a+h,b)-f(a,b)}{h},\qquad f_y(a,b)=\lim_{k\to0}\frac{f(a,b+k)-f(a,b)}{k}
\]
Hesap: diğer değişken(ler) sabit. Geometri: $f_x$, $f_y$ yüzeyin $y=b$ ve $x=a$ düzlemlerindeki izlerinin eğimleri.

\textbf{İkinci mertebe:} $f_{xx},\ f_{xy}=(f_x)_y=\dfrac{\partial^2f}{\partial y\,\partial x},\ f_{yx},\ f_{yy}$.\\
\textbf{Schwarz/Clairaut:} $f_{xy}$, $f_{yx}$ sürekli $\Rightarrow f_{xy}=f_{yx}$.

\textbf{Diferansiyellenebilirlik:}
$\Delta z=f_x\Delta x+f_y\Delta y+\varepsilon_1\Delta x+\varepsilon_2\Delta y$, \ $\varepsilon_1,\varepsilon_2\to0$.\\
$f_x,f_y$ sürekli $\Rightarrow$ diferansiyellenebilir $\Rightarrow$ sürekli (tersleri genelde yanlış).

\begin{center}
\begin{tabular}{l l}
\toprule
Kavram & Formül \\
\midrule
Teğet düzlem & $z-z_0=f_x(x_0,y_0)(x-x_0)+f_y(x_0,y_0)(y-y_0)$ \\
Normal doğru & $(x_0,y_0,z_0)+t\,\langle f_x,f_y,-1\rangle$ \\
Lineerleştirme & $L(x,y)=f(x_0,y_0)+f_x(x_0,y_0)(x-x_0)+f_y(x_0,y_0)(y-y_0)$ \\
Hata sınırı & $|E|\le\frac M2\big(|x-x_0|+|y-y_0|\big)^2$, \ $M\ge|f_{xx}|,|f_{xy}|,|f_{yy}|$ ($B$'de) \\
Tam diferansiyel & $dz=f_x\,dx+f_y\,dy$, \quad $dw=f_x\,dx+f_y\,dy+f_z\,dz$ \\
İkinci diferansiyel & $d^2z=z_{xx}dx^2+2z_{xy}dx\,dy+z_{yy}dy^2$ \\
$n$-inci diferansiyel & $d^nz=\left(\frac{\partial}{\partial x}dx+\frac{\partial}{\partial y}dy\right)^{(n)}z$ \\
Yaklaşık hesap & $\Delta z\approx dz$, \quad $f(x_0+dx,y_0+dy)\approx f(x_0,y_0)+dz$ \\
Bağıl / yüzde & $dz/z$, \quad $100\cdot dz/z$ (\%) \\
Kuvvet çarpımı & $z=Cx^ay^b\Rightarrow dz/z=a\,dx/x+b\,dy/y$ \\
\bottomrule
\end{tabular}
\end{center}

\textbf{Üç değişken:} $L=f(P_0)+f_x\Delta x+f_y\Delta y+f_z\Delta z$, \
$|E|\le\frac M2(|\Delta x|+|\Delta y|+|\Delta z|)^2$, \
$d^2w=\sum w_{ii}\,dx_i^2+2\sum_{i<j}w_{ij}\,dx_i\,dx_j$.
\end{ozet}
